
%%% HUJ







% The clean fix would have been to use ntHash~\cite{nthash_2016}, the
% rolling hash recommended by the SIMD-Minimizers
% paper~\cite{groot_koerkamp_simd_minimizers_2025}.  ntHash is a
% sequence of \texttt{rotate\,+\,XOR\,+\,table-lookup} operations
% where the table is indexed by the 2-bit nucleotide rank, so on NEON
% it can be implemented as a $4$-entry \texttt{tbl.16b} lookup per
% character.  This produces a uniform hash distribution across the
% $k$-mer space, which removes the AT-bias.

% V14 implements canonical ntHash inside the V13 dual-stream chassis.
% The implementation is correct -- it preserves the warmup-cancellation
% invariant by precomputing constants that subtract the
% $c_\text{old} = 0$-padding contribution after $k$ rolling steps --
% and the integration suite passes.

% What V14 \emph{didn't} buy was wall-clock:

% \begin{itemize}
% \item \texttt{mm\_sketch\_simd} cost \emph{increased} by 2\,ms
%   (the table lookup adds a NEON \texttt{tbl.16b}, two rotates, and
%   some XORs per stream per BP).
% \item \texttt{collect\_seed\_hits} dropped by \emph{0.4\,ms}, not
%   the predicted 5\,ms.
% \item Net: V14 is statistically tied with V13 at LONG and
%   VERY\_LONG, and slightly slower at XL.
% \end{itemize}

% The lesson is that the AT-bias hypothesis was wrong.  The
% \texttt{collect\_seed\_hits} regression is dominated not by the
% minimizer hash distribution but by the \texttt{std::unordered\_map}
% lookup overhead: bioseq's target index uses STL's default hash
% table, while the C minimap2 path uses an open-addressing
% \texttt{khash\_t<uint64\_t>}.  The next-most-impactful single edit
% for closing Goal~2 -- replacing \texttt{std::unordered\_map} with a
% flat hash on a uint32 $\to$ vector layout -- is in the post-V14
% future-work list of the engineering report and is not part of this
% chapter.

% \subsection{V15 -- dual-batch chaining}
% \label{sec:v15-chain-v2}

% The dual-stream interleaving idea that delivered for the minimizer
% loop (\S\ref{sec:v11-v12}) also worked beautifully for an aggressive
% unrolling of the SIMD chaining inner loop. Unfortunately, the
% chaining loop is not register-bound and 4-pipe-issue-limited;
% \emph{it is load-store-bound}.  V15 doubles the SoA loads per inner
% iteration from 4 to 8, which queues up on the M1's 4 load-store
% pipes and stalls the issue queue.  Net wall-clock effect:

% \begin{itemize}
% \item \texttt{mg\_lchain\_simd} regresses by 5--13\,\% at 903\,kbp.
% \item Total pipeline regresses by 7\,\% versus V13 chaining + V13 minimizer.
% \end{itemize}

% The lesson is that the unrolling-with-multiple-reduction-streams
% pattern only helps register-bound, issue-width-limited kernels.
% Apply it only after confirming the target kernel is issue-width
% limited, not load-store limited.  V15 stays in the codebase as a
% documented negative result.

% \subsection{Goal~2-style \texttt{prefix\_kmer} drop}

% A future-work item from V12 proposed dropping the parallel
% \texttt{ring\_buf\_hs} buffer and the \texttt{prefix\_kmer} state
% register by packing the strand bit into bit~16 of \texttt{elem}.
% The estimated saving was $\sim 3$ ops/stream/iter and 22 vec\_t of
% register-pressure relief.  We implemented and benchmarked.  Result:

% \begin{verbatim}
%                                   V12 (pre)   V12 (post-Goal-2)   delta
% mm_sketch_simd                    9.0 ms      9.5 ms              +0.5 ms (matches estimate)
% target_index_build                24.7 ms     4.9 ms              -19.8 ms (~85% fewer keys)
% collect_seed_hits                 12.3 ms     842.0 ms            +830 ms (70x regression)
% mg_lchain_simd                    8.9 ms      299.9 ms            +291 ms (34x regression)
% mm_align_skeleton                 87.2 ms    4811.0 ms            +4724 ms (56x regression)
% \end{verbatim}

% The kmer space stored in \texttt{elem} after the trick was only 14
% bits -- not the 16 bits the V12 future-work entry assumed -- so the
% \texttt{target\_index} hash buckets collapsed to $\sim 16\,000$
% distinct keys for $168\,000$ query minimizers.  Each query had
% $\sim 10$ collisions per bucket, the chain-anchor count exploded,
% and the alignment skeleton had to walk through all of them.  The
% pipeline is $36\times$ slower as a result.

% The lesson is that hash-precision losses in a sketching scheme
% \emph{compound} catastrophically through the downstream pipeline.
% Even an apparently small precision loss must be validated against
% the full pipeline's wall-clock, not just the sketch phase's.

% \section{Round 5 -- the divergence-estimator fix (V13)}
% \label{sec:v13-est-err}

% The single biggest \emph{non-minimizer} contribution to closing
% Goal~2 is the V13 rewrite of
% \texttt{estimate\_divergence\_from\_seeds} in
% \path{bioseq/backend/sequences/minimap_port.cppm}, which dropped
% \texttt{mm\_est\_err} from V12-pre's 7.88\,ms to V13's 0.85\,ms --
% a $9\times$ speedup, removing one of the two large downstream
% regressions in the SIMD pipeline.

% \subsection{Why \texttt{mm\_est\_err} mattered}

% For each chain, minimap2 estimates the per-base sequence divergence
% from a Poisson model on the minimizer-set
% intersection~\cite{minimap2_paper}:
% \[
% \hat{\epsilon} = \frac{1}{k} \log \frac{n}{m}
% \]
% where $n$ is the number of query minimizers in the chain's reference
% span and $m$ is the number of those that are actually seeded
% (i.e.\ that appear as anchors in the chain).  Computing $n$ requires
% \emph{knowing which positions are query minimizers}, which is a
% membership test against the query's minimizer-position set.

% The pre-V13 implementation:

% \begin{enumerate}
% \item Built a sorted array \texttt{minimizer\_positions[]} of every
%   query minimizer's position via \texttt{std::sort}.  The V12
%   dual-stream collector emits in lane-interleaved order, so the
%   array is \emph{not} sorted on output, and \texttt{std::sort}
%   sees worst-case input.
% \item For each chain anchor, did a \texttt{std::lower\_bound} into
%   the sorted array to test membership.
% \end{enumerate}

% The dominant cost was the \texttt{std::sort} on 168\,k entries:
% $\sim 5$--$6$\,ms per pair.  The \texttt{lower\_bound} loop added
% another $\sim 1$--$2$\,ms.

% \subsection{The sort-free prefix-count rewrite}

% V13 replaces both with a single linear-pass prefix-sum:

% \[
% \textit{prefix\_count}[i] = \big|\{ \, p \in \textit{minimizer\_positions} \mid p < i \,\}\big|
% \]

% Built in $O(N + L)$ where $N$ is the minimizer count and $L$ is the
% query length.  After construction:

% \begin{itemize}
% \item $\textit{is\_minimizer}(p) = (\textit{prefix\_count}[p+1] > \textit{prefix\_count}[p])$ -- $O(1)$.
% \item $\textit{range\_count}(a, b) = \textit{prefix\_count}[b+1] - \textit{prefix\_count}[a]$ -- $O(1)$.
% \end{itemize}

% The per-region inner loop is now a single $O(\textit{seed\_count})$
% walk via $O(1)$ probes.  No sort, no binary search.

% The fix landed for \emph{all} configs (scalar pipeline included) so
% even the scalar path got a small speedup: scalar
% \texttt{mm\_est\_err} dropped from 1.21\,ms to 0.85\,ms, V12-SIMD
% \texttt{mm\_est\_err} dropped from 7.88\,ms to 0.85\,ms.  The
% SIMD-minus-scalar regression is gone.

% There was also a correctness bug exposed by the rewrite: the V10
% ``reverse-chain fix'' was wrong, because for both forward and
% reverse chains the lambda \texttt{compute\_forward\_pos} was
% already producing monotonically-increasing forward positions in
% chain order. We removed the reverse-chain split entirely.

% \subsection{Net pipeline impact}

% Figure~\ref{fig:phase-stack-v13} stacks the per-phase cost at
% 903\,kbp for V12-pre versus V13.  The two large gains are
% \texttt{mm\_est\_err} (from 7.88\,ms to 0.85\,ms) and the V12
% \texttt{mm\_sketch} stack (from 14.99\,ms to 9.23\,ms, already
% covered in \S\ref{sec:v11-v12}).

% \begin{figure}[ht]
%   \centering
%   \begin{CacheMeCode}{python, custom include command={\input{\robExtAddCachePathAndName{\robExtFinalHash.pgf}}}, set placeholder eval={__LINEWIDTH__}{\lenToCmNoUnit[in]{\linewidth}}}
%     import matplotlib.pyplot as plt
%     import matplotlib
%     import numpy as np
%     matplotlib.use("pgf")
%     matplotlib.rcParams.update({
%       "font.family": "serif",
%       "font.serif": [],
%       "text.usetex": True,
%     })
%     fig, ax = plt.subplots()
%     fig.set_size_inches(__LINEWIDTH__, 0.55*__LINEWIDTH__, forward=True)
%     configs = ['scalar', 'V6 SIMD', 'V8 two-stack', 'V12 v3 (pre-V13)', 'V12 v3 (post-V13)']
%     phases = [
%         ('mm\\_sketch + mm\\_sketch\\_simd', [12.50, 38.82, 32.43, 14.99, 14.52], '#9999cc'),
%         ('target\\_index\\_build',         [16.0,  14.3,  14.6,  26.38, 23.83], '#1f77b4'),
%         ('collect\\_seed\\_hits',          [6.5,   11.0,  11.3,  12.32, 12.0],  '#2ca02c'),
%         ('mg\\_lchain\\_dp/simd',          [16.0,  8.9,   8.9,   8.88,  8.78],  '#ff7f0e'),
%         ('mm\\_est\\_err',                 [0.7,   7.2,   7.1,   7.88,  0.85],  '#d62728'),
%         ('mm\\_align\\_skeleton',          [87.5,  86.4,  86.5,  87.25, 85.42], '#888888'),
%     ]
%     bottoms = np.zeros(len(configs))
%     for name, vals, color in phases:
%         ax.bar(configs, vals, bottom=bottoms, label=name, color=color)
%         bottoms = bottoms + np.array(vals)
%     for i, b in enumerate(bottoms):
%         ax.text(i, b + 1.5, f'{b:.0f}', ha='center', fontsize=7)
%     ax.set_ylabel('Per-pair phase cost (ms) at 903\\,kbp', fontsize=8)
%     ax.set_title('Per-phase pipeline composition (903\\,kbp, mean of 3 runs)', fontsize=9)
%     ax.tick_params(axis='both', labelsize=7)
%     ax.legend(fontsize=6, frameon=False, loc='upper right', ncol=2)
%     ax.grid(True, alpha=0.3, axis='y')
%     plt.xticks(rotation=15, ha='right', fontsize=7)
%     print(get_filename_from_extension(".pgf"))
%     plt.savefig(get_filename_from_extension(".pgf"), bbox_inches="tight")
%   \end{CacheMeCode}
%   \caption{Per-phase composition of the per-pair pipeline at
%     903\,kbp for the headline configurations.  ``V12 v3 (pre-V13)''
%     is the V12 collector with the pre-V13 \texttt{mm\_est\_err}
%     code; ``V12 v3 (post-V13)'' is the same minimizer collector with
%     the V13 sort-free prefix-count rewrite.  The pre/post pair lets
%     the reader visually separate the V12 (minimizer-side) wins from
%     the V13 (\texttt{mm\_est\_err}) wins.  The \texttt{mm\_align\_skeleton}
%     bar is identical across all configs, accounting for $\sim 50\,\%$
%     of total wall-clock, and is not addressed by any of the work in
%     this chapter.}
%   \label{fig:phase-stack-v13}
% \end{figure}

% \section{Final results}
% \label{sec:final-results}

% The combination of the V12 minimizer collector and the V13
% \texttt{mm\_est\_err} rewrite is the production state at the time of
% writing.  Across the five sequence-length buckets the bench
% exercises (LONG, VERY\_LONG, XL), V13 is the fastest of the SIMD
% configurations on every bucket and is the fastest config
% \emph{overall} on the LONG sub-buckets (47\,kbp, 49\,kbp, 52\,kbp).
% At VERY\_LONG (903\,kbp) it is 8.8\,\% faster than scalar -- just
% short of the 10\,\% Goal-2 target.  At XL (5\,Mbp) it is 4.9\,\%
% faster than scalar.

% \begin{figure}[ht]
%   \centering
%   \begin{CacheMeCode}{python, custom include command={\input{\robExtAddCachePathAndName{\robExtFinalHash.pgf}}}, set placeholder eval={__LINEWIDTH__}{\lenToCmNoUnit[in]{\linewidth}}}
%     import matplotlib.pyplot as plt
%     import matplotlib
%     import numpy as np
%     matplotlib.use("pgf")
%     matplotlib.rcParams.update({
%       "font.family": "serif",
%       "font.serif": [],
%       "text.usetex": True,
%     })
%     fig, ax = plt.subplots()
%     fig.set_size_inches(__LINEWIDTH__, 0.50*__LINEWIDTH__, forward=True)
%     buckets = ['47 kbp', '49 kbp', '52 kbp', '903 kbp', '5.0 Mbp']
%     delta_v6   = [(15.21 - 14.56) / 14.56 * 100,
%                   (15.90 - 14.94) / 14.94 * 100,
%                   (16.05 - 15.09) / 15.09 * 100,
%                   (201.68 - 188.93) / 188.93 * 100,
%                   (1207.99 - 1171.0) / 1171.0 * 100]
%     delta_v8   = [(14.64 - 14.56) / 14.56 * 100,
%                   (15.07 - 14.94) / 14.94 * 100,
%                   (15.24 - 15.09) / 15.09 * 100,
%                   (183.77 - 188.93) / 188.93 * 100,
%                   (1131.87 - 1171.0) / 1171.0 * 100]
%     delta_v13  = [(13.78 - 14.56) / 14.56 * 100,
%                   (14.02 - 14.94) / 14.94 * 100,
%                   (13.82 - 15.09) / 15.09 * 100,
%                   (172.38 - 188.93) / 188.93 * 100,
%                   (1113.3 - 1171.0) / 1171.0 * 100]
%     x = np.arange(len(buckets))
%     width = 0.27
%     ax.bar(x - width, delta_v6,  width, label='V6 SIMD (rescan)', color='#1f77b4')
%     ax.bar(x,         delta_v8,  width, label='V8 two-stack',     color='#2ca02c')
%     ax.bar(x + width, delta_v13, width, label='V13 (production)', color='#d62728')
%     ax.axhline(-10, linestyle='--', linewidth=1, color='#ff0000',
%                 label='-10\\% Goal-2 target')
%     ax.axhline(0, linewidth=0.6, color='#000000')
%     ax.set_xticks(x)
%     ax.set_xticklabels(buckets)
%     ax.set_ylabel('Wall-clock delta vs scalar (\\%, negative = faster)', fontsize=8)
%     ax.set_title('SIMD vs scalar speedup across iterations', fontsize=9)
%     ax.tick_params(axis='both', labelsize=7)
%     ax.legend(fontsize=7, frameon=False, loc='upper right', ncol=2)
%     ax.grid(True, alpha=0.3, axis='y')
%     print(get_filename_from_extension(".pgf"))
%     plt.savefig(get_filename_from_extension(".pgf"), bbox_inches="tight")
%   \end{CacheMeCode}
%   \caption{Wall-clock delta versus scalar for the three headline
%     SIMD configurations, expressed as percent (negative = faster).
%     The dashed line is the $-10\,\%$ Goal-2 target. V13 approaches
%     the target across the LONG / VERY\_LONG buckets (between
%     $-8.4\,\%$ and $-8.8\,\%$ on the plotted three-run means) and
%     is within 5\,\% of the target on XL.}
%   \label{fig:speedup-summary}
% \end{figure}

% The summary table:

% \begin{table}[ht]
%   \centering
%   \small
%   \begin{tabular}{lrrrr}
%     \toprule
%     Length     & scalar (ms) & V13 (ms) & V13 vs scalar & V13 vs V6 SIMD \\
%     \midrule
%     47\,kbp   &  14.56  & 13.78 & $-5.4\,\%$ & $-9.4\,\%$ \\
%     49\,kbp   &  14.94  & 14.02 & $-6.2\,\%$ & $-11.8\,\%$ \\
%     52\,kbp   &  15.09  & 13.82 & $-8.4\,\%$ & $-13.9\,\%$ \\
%     903\,kbp  & 188.93  & 172.38 & $-8.8\,\%$ & $-14.5\,\%$ \\
%     5.0\,Mbp  & 1171.0  & 1113.3 & $-4.9\,\%$ & $-7.8\,\%$ \\
%     \bottomrule
%   \end{tabular}
%   \caption{Final V13-state speedup numbers across the five
%     sequence-length buckets, expressed as wall-clock per pair
%     (mean of 3 runs from
%     \texttt{[quick-long-sequence-benchmark]}, bench-opt
%     \texttt{-O3 -DNDEBUG}, CCS error rate, seed
%     \texttt{0xBEEFCAFE12345678}). The $-10\,\%$ Goal-2 target is
%     narrowly missed at 903\,kbp and 5.0\,Mbp; on the LONG sub-buckets
%     (47, 49, 52\,kbp) it is met or comfortably approached.}
%   \label{tab:final-numbers}
% \end{table}

% $$$ SUPEREND

% \begin{table}[H]
%   \centering
%   \small
%   \begin{tabularx}{\linewidth}{p{.20\linewidth}XX}
%   \toprule
%   Method & What happens when one candidate enters & Why it matters here \\
%   \midrule
%   Saved minimum plus rescan & Compare once; scan all $w$ slots only when the
%   saved record leaves & Excellent scalar common case; one expiring SIMD lane
%   starts shared work \\
%   Monotone deque & Remove expired records at the front and larger records at the
%   back & The number of removals depends on the hash values \\
%   Prefix and suffix minima & Rebuild a fixed block summary every $w$ steps and
%   combine two saved values per crossing window & The instruction schedule
%   depends on the step number, not on which lane contains the minimum \\
%   \bottomrule
%   \end{tabularx}
%   \caption{Three ways to maintain a sliding minimum. The next section uses the
%   third method because every lane performs the same fixed steps.}
%   \label{tab:rewrite-sliding-minimum-methods}
%   \end{table}

%   \paragraph{Is the specialized kernel different from the generic kernel?}
%   Yes in implementation, but not in minimizer semantics: both calculate
%   chronological adjacent windows, choose the rightmost equal minimum, and use
%   the same transition and tie-output rules.
  
%   \begin{enumerate}
%     \item \textbf{Batching.}  A generic full call is one $L$-window group.
%           The fast call stages eight candidates and runs two sequential
%           four-window groups; eight is fixed in the code, not derived from
%           $k$ or $w$.
%     \item \textbf{Construction.}  The generic kernel constructs $L$ exact
%           $k$-mers in $k$ shifted-vector steps.  The fast kernel constructs
%           eight exact K15 records with fixed SWAR operations in a 24-byte
%           scratch area: 14 retained ranks, eight new ranks, and two safety
%           bytes.
%     \item \textbf{Reduction.}  The generic group reduces
%           $(w-1)+L$ records in $w$ steps.  Each fast group reduces
%           $9+4=13$ records in ten steps, updates the tail, and feeds the next
%           group; a full call therefore performs 20 reduction steps.
%     \item \textbf{Applicability.}  The fast route is restricted to its
%           K15/W10 DNA4 packed-u32 layout and vector target.  Other admitted
%           parameters, and ambiguity-aware DNA5 input, use the generic route.
%   \end{enumerate}
  
%   \subsection{Algorithmic differences: fixed and probability-dependent work}
  
%   Let $N$ be the number of useful candidate positions and let $L$ be the one
%   lane-count symbol used for all three methods.  Thus
%   $\lceil N/L\rceil$ is the number of SIMD groups.  A1 and A2 place separated
%   sequence streams in those lanes, whereas A3 places adjacent windows in them.
  
%   Use the requested two-event model.  Let $p$ be the per-lane probability that
%   the saved minimum expires, and let $r>p$ be the per-lane probability that a
%   window contains a repeated minimum hash.  To obtain a closed form from only
%   $p$ and $r$, assume these events are independent within a lane and assume
%   independent lanes.  This is a deliberately simplified estimate because
%   adjacent A3 lanes share $w-1$ candidates.  Define
%   \[
%     \pi_p=1-(1-p)^L,
%     \qquad
%     \pi_{pr}=1-(1-pr)^L .
%   \]
%   The first quantity estimates the probability that some lane expires.  The
%   second estimates the probability that some lane both expires and contains a
%   repeated minimum hash.  A raw-distinct output tie is a subset of the latter
%   event, so terms using $\pi_{pr}$ below are upper bounds.  Notice that $r>p$
%   does not make the tie branch more frequent than expiry: $pr\leq p$, hence
%   $\pi_{pr}\leq\pi_p$.  One-time startup and final-tail effects are included in
%   the additive $O(k+w)$ term.
  
%   \begin{table}[H]
%   \centering
%   \scriptsize
%   \begin{tabularx}{\linewidth}{p{.11\linewidth}p{.25\linewidth}
%                                p{.29\linewidth}X}
%   \toprule
%   Method & Fixed-schedule work & Probability-dependent work
%          & Expected and worst-case work \\
%   \midrule
%   A1: naive SIMD rescan &
%   One constant-cost advance per group after $O(k+w)$ setup. &
%   With probability $\pi_p$, the shared slow path performs one $w$-slot minimum
%   rescan and one $w$-slot equal-minimum scan.  The value of $r$ changes output,
%   not this schedule. &
%   \[
%    \begin{aligned}
%    \mathbb{E}[T_{\mathrm{A1}}]=O\biggl(&k+w\\[-2pt]
%      &+\left\lceil\frac NL\right\rceil(1+2w\pi_p)\biggr).
%    \end{aligned}
%   \]
%   Worst case:
%   $O\!\left(k+w+\lceil N/L\rceil w\right)$. \\
%   \addlinespace
%   A2: rolling ntHash and prefix--suffix minima &
%   Rolling hash update and two aggregate merges per group; one deterministic
%   $(w-1)$-merge suffix rebuild every $w$ groups. &
%   None.  Expiry and ties do not change the instruction schedule. &
%   Minimum maintenance performs
%   \[
%     \begin{aligned}
%     &2\left\lceil\frac NL\right\rceil\\[-2pt]
%     &\quad +(w-1)\left\lfloor
%         \frac{\lceil N/L\rceil}{w}\right\rfloor
%     \end{aligned}
%   \]
%   merges.  Expected and worst case:
%   $O(k+w+N/L)$. \\
%   \addlinespace
%   A3: generic chronological reduction &
%   Every group performs $k$ shifted exact-$k$-mer construction steps and $w$
%   shifted adjacent-window reductions:
%   \[
%     \begin{aligned}
%     \Theta\biggl(&k+w\\[-2pt]
%       &+\left\lceil\frac NL\right\rceil(k+w)\biggr).
%     \end{aligned}
%   \]
%   Leading term: $\Theta(N(k+w)/L)$. &
%   An expiry-gated group adds a $w$-slot tie-verification pass.  A confirmed tie
%   adds an event construction, transpose, ordering, and compression path costing
%   $O(L(w+L))$ per affected group.  Neither path recomputes a minimum. &
%   \[
%    \begin{aligned}
%    \mathbb{E}[T_{\mathrm{A3}}]=O\biggl(
%      &k+w\\[-2pt]
%      &+\frac NL[k+w(1+\pi_p)]\\[-2pt]
%      &+N(w+L)\pi_{pr}\biggr).
%    \end{aligned}
%   \]
%   Worst case:
%   $O\!\left(k+N[k/L+w+L]\right)$. \\
%   \bottomrule
%   \end{tabularx}
%   \caption{Fixed and conditional source-level work using only $N,L,w,k,p,r$.
%   The table counts vector-loop and tie-layout work, not machine instructions.}
%   \label{tab:rewrite-fixed-variable-work}
%   \end{table}
  
%   For the fixed K15/W10 specialization, substitute $k=15$, $w=10$, and $L=4$.
%   A full eight-candidate call is simply two consecutive four-lane groups, so
%   the sustained stream still contains $\lceil N/L\rceil$ groups.  The numeral
%   eight never enters the asymptotic expression.  Under the two-probability
%   model above, the following is a conservative upper estimate:
%   \[
%     O\!\left(
%       k+w+\left\lceil\frac{N}{2L}\right\rceil
%       +\left\lceil\frac NL\right\rceil
%       \left[
%         w(1+\pi_p+2\pi_{pr})+L(w+L)\pi_{pr}
%       \right]
%     \right).
%   \]
%   The first ceiling counts one fixed packed K15 construction per pair of
%   four-lane groups.  The bracket contains the always-run minimum reduction, the
%   expiry hash prefilter, the tie raw-verification and event-enumeration passes,
%   and the rare event layout.  With $k=15,w=10,L=4$ all of these are constants,
%   so expected and worst-case time are both $\Theta(N)$.
  
%   For one full fast call, the minimum loop itself always performs
%   \[
%     \underbrace{10}_{\text{first group}}
%     +\underbrace{10}_{\text{second group}}
%     =20
%   \]
%   shifted hash comparisons.  Its expected count of ten-slot reduction or
%   verification loops is bounded by
%   \[
%     20(1+\pi_p+2\pi_{pr}),
%     \qquad
%     \pi_p=1-(1-p)^4,
%     \qquad
%     \pi_{pr}=1-(1-pr)^4.
%   \]
%   The two $\pi_{pr}$ terms conservatively charge both raw verification and
%   event enumeration whenever the repeated-hash gate succeeds; the source can
%   reject the event before the latter pass.  The fixed twenty reduction steps
%   load 20 u32 hash vectors, 40 u64 raw-record
%   vectors, and 40 u64 ordinal vectors.  Expiry adds no minimum-reduction load;
%   it can only trigger the separate verification work shown above.
  
%   Algorithmically, A2 has the best guaranteed scaling in $k$ and $w$: after
%   initialization it is $O(N/L)$ and its schedule is independent of $p$ and $r$.
%   A1 can execute fewer minimum operations when expiry is extremely rare, but
%   its worst case grows as $Nw/L$.  Generic A3 deliberately pays
%   $N(k+w)/L$ fixed work to calculate complete adjacent windows.  These orders do
%   not predict fixed-$k$, fixed-$w$ wall-clock speed, where the specialized
%   packing and the different metadata operations can dominate.
  

% \begin{figure}[H]
%   \centering
%   \begin{tikzpicture}[
%     x=.56cm,
%     y=.56cm,
%     font=\footnotesize,
%     >=Latex,
%     base old/.style={draw=black!70,fill=rewriteGray!18,
%       minimum width=.36cm,minimum height=.42cm,inner sep=0pt},
%     base new/.style={draw=rewriteBlue!90!black,fill=rewriteBlue!14,
%       minimum width=.36cm,minimum height=.42cm,inner sep=0pt},
%     base next/.style={draw=rewriteBlue!90!black,densely dashed,
%       fill=rewriteBlue!7,minimum width=.36cm,minimum height=.42cm,inner sep=0pt},
%     unused/.style={draw=rewriteGray!75,dashed,fill=rewriteGray!6,
%       text=rewriteGray!85!black,minimum width=.36cm,
%       minimum height=.42cm,inner sep=0pt},
%     slot/.style={draw=black!65,fill=white,
%       minimum width=.36cm,minimum height=.34cm,inner sep=0pt},
%     picked/.style={draw=rewriteOrange!90!black,very thick,
%       fill=rewriteOrange!20,minimum width=.36cm,
%       minimum height=.34cm,inner sep=0pt},
%     tied/.style={draw=rewriteOrange!90!black,densely dashed,
%       fill=rewriteOrange!7,minimum width=.36cm,
%       minimum height=.34cm,inner sep=0pt},
%     answer/.style={draw=rewriteOrange!90!black,very thick,
%       fill=rewriteOrange!14,rounded corners,
%       minimum width=1.45cm,minimum height=.40cm,inner sep=2pt},
%     flow/.style={-{Latex[length=1.8mm]},thick}
%   ]
  
%   % (a) Generic dimensions and the only partial call.
%   \node[anchor=west,font=\bfseries] at (0,21.35)
%     {(a) One ordinary generic DNA4 iteration processes exactly $L$ candidates};
  
%   \node[anchor=east] at (1.90,20.45) {base ranks};
%   \node[base old,minimum width=4.20cm] (goldbase) at (5.95,20.45)
%     {$k-1$ retained};
%   \node[base new,minimum width=2.25cm] (gnewbase) at (11.71,20.45)
%     {$L$ new};
%   \draw[flow] (gnewbase.east)--(15.10,20.45);
%   \node[anchor=west] at (15.25,20.45)
%     {$k$ shifted updates $\longrightarrow L$ hashes};
%   \draw[decorate,decoration={brace,mirror,amplitude=3pt}]
%     (2.18,20.02)--(13.72,20.02)
%     node[midway,below=4pt] {$(k-1)+L$ live ranks};
  
%   \node[anchor=east] at (1.90,18.70) {candidate row};
%   \node[base old,minimum width=4.20cm] (goldcand) at (5.95,18.70)
%     {$w-1$ retained};
%   \node[base new,minimum width=2.25cm] (gnewcand) at (11.71,18.70)
%     {$L$ new};
%   \draw[flow] (gnewcand.east)--(15.10,18.70);
%   \node[anchor=west] at (15.25,18.70)
%     {$w$ shifted reductions $\longrightarrow L$ minima};
%   \draw[decorate,decoration={brace,mirror,amplitude=3pt}]
%     (2.18,18.27)--(13.72,18.27)
%     node[midway,below=4pt] {$(w-1)+L$ candidate records};
  
%   \node[anchor=east] at (1.90,16.85) {ordinary};
%   \node[base new,minimum width=2.25cm] (fullpending) at (4.15,16.85)
%     {$L$ pending};
%   \draw[flow] (fullpending.east)--(7.00,16.85);
%   \node[anchor=west] at (7.15,16.85) {run one full group};
  
%   \node[anchor=east] at (12.50,16.85) {end of input only};
%   \node[base new,minimum width=1.25cm] (tailactive) at (14.10,16.85)
%     {$q$};
%   \node[unused,minimum width=1.75cm] (tailidle) at (16.80,16.85)
%     {$L-q$};
%   \draw[flow] (tailidle.east)--(19.35,16.85);
%   \node[anchor=west] at (19.50,16.85) {one masked group};
  
%   % (b) Fixed K15/W10 staging and the two chronological groups.
%   \node[anchor=west,font=\bfseries] at (0,14.25)
%     {(b) The fixed $k=15,w=10$ fast call stages eight new candidates};
  
%   \node[anchor=east] at (1.85,13.35) {base scratch};
%   \foreach \i in {0,...,13}
%     \node[base old,minimum width=.27cm] at (2.35+.52*\i,13.35) {};
%   \foreach \i in {14,...,17}
%     \node[base new,minimum width=.27cm] at (2.35+.52*\i,13.35) {};
%   \foreach \i in {18,...,21}
%     \node[base next,minimum width=.27cm] at (2.35+.52*\i,13.35) {};
%   \foreach \i in {22,23}
%     \node[unused,minimum width=.27cm] at (2.35+.52*\i,13.35) {$0$};
  
%   \draw[decorate,decoration={brace,mirror,amplitude=3pt}]
%     (2.08,12.93)--(9.38,12.93)
%     node[midway,below=3pt] {$14$ retained};
%   \draw[decorate,decoration={brace,mirror,amplitude=3pt}]
%     (9.42,12.93)--(11.46,12.93)
%     node[midway,below=3pt] {first $4$};
%   \draw[decorate,decoration={brace,mirror,amplitude=3pt}]
%     (11.50,12.93)--(13.54,12.93)
%     node[midway,below=3pt] {next $4$};
%   \draw[decorate,decoration={brace,mirror,amplitude=3pt}]
%     (13.58,12.93)--(14.58,12.93)
%     node[midway,below=3pt] {$2$};
  
%   \node[anchor=west] at (16.00,13.35)
%     {$14+8=22$ live ranks};
%   \node[anchor=west,text=rewriteGray!90!black] at (16.00,12.70)
%     {fixed scratch: $24$ bytes};
  
%   % (c) One complete group reduction.
%   \node[anchor=west,font=\bfseries] at (0,10.20)
%     {(c) Group 0: four adjacent windows are reduced in the same ten SIMD steps};
  
%   \node[anchor=east] at (1.85,9.35) {row};
%   \foreach \i/\v/\sty in {
%     0/8/base old,1/6/base old,2/5/base old,3/9/base old,
%     4/4/base old,5/7/base old,6/4/base old,7/8/base old,8/3/base old,
%     9/6/base new,10/3/base new,11/5/base new,12/2/base new}
%     \node[\sty] at (2.35+.72*\i,9.35) {\v};
%   \foreach \i in {0,...,12}
%     \node[font=\tiny,anchor=north] at (2.35+.72*\i,9.00) {$D_{\i}$};
  
%   \foreach \ell/\first/\winner/\value/\result in {
%     0/0/8/3/$3@D_8$,
%     1/1/10/3/$3@D_{10}$,
%     2/2/10/3/$3@D_{10}$,
%     3/3/12/2/$2@D_{12}$} {
%     \pgfmathsetmacro{\yy}{8.15-.72*\ell}
%     \node[anchor=east] at (1.85,\yy) {$W_{\ell}$};
%     \foreach \j in {0,...,9} {
%       \pgfmathtruncatemacro{\idx}{\first+\j}
%       \pgfmathparse{int(\idx==\winner)}
%       \ifnum\pgfmathresult=1
%         \node[picked] at (2.35+.72*\idx,\yy) {\value};
%       \else
%         \node[slot] at (2.35+.72*\idx,\yy) {};
%       \fi
%     }
%     \node[answer] at (12.95,\yy) {\result};
%   }
%   \node[tied] at (2.35+.72*8,7.43) {3};
%   \node[tied] at (2.35+.72*8,6.71) {3};
  
%   \draw[flow,rewriteOrange!90!black]
%     (2.35,5.30)--(11.65,5.30)
%     node[midway,below=3pt]
%       {$r=0,\ldots,w-1$: $V_r=(D_r,D_{r+1},D_{r+2},D_{r+3})$};
%   \node[anchor=west,text=rewriteOrange!90!black] at (14.70,6.55)
%     {$10$ shifted reductions};
%   \node[anchor=west,text=rewriteGray!90!black] at (14.70,5.90)
%     {dashed: earlier equal minimum};
  
%   % (d) State update and the second group.
%   \node[anchor=west,font=\bfseries] at (0,4.35)
%     {(d) Advance four positions, retain the overlap, and repeat};
  
%   \node[anchor=east] at (1.85,3.55) {group 1 row};
%   \foreach \i/\v/\sty in {
%     0/4/base old,1/7/base old,2/4/base old,3/8/base old,4/3/base old,
%     5/6/base old,6/3/base old,7/5/base old,8/2/base old,
%     9/7/base next,10/1/base next,11/9/base next,12/1/base next}
%     \node[\sty] at (2.35+.72*\i,3.55) {\v};
%   \draw[decorate,decoration={brace,mirror,amplitude=3pt}]
%     (2.03,3.13)--(8.43,3.13)
%     node[midway,below=3pt] {new tail $D_4,\ldots,D_{12}$};
%   \draw[decorate,decoration={brace,mirror,amplitude=3pt}]
%     (8.51,3.13)--(11.31,3.13)
%     node[midway,below=3pt] {$D_{13},\ldots,D_{16}$};
  
%   \node[anchor=east] at (15.55,3.55) {saved current};
%   \node[answer] at (17.25,3.55) {$2@D_{12}$};
  
%   \draw[flow] (1.00,3.15)--(1.00,1.95);
%   \node[anchor=west,text=rewriteOrange!90!black] at (14.50,2.55)
%     {same ten reductions};
  
%   \node[anchor=east] at (1.45,1.55) {minima};
%   \foreach \i/\v in {
%     0/$2@D_{12}$,1/$1@D_{14}$,2/$1@D_{14}$,3/$1@D_{16}$}
%     \node[answer,minimum width=1.50cm] at (3.10+3.10*\i,1.55) {\v};
%   \node[anchor=west] at (14.50,1.85)
%     {final tail $D_8,\ldots,D_{16}$};
%   \node[anchor=west] at (14.50,1.20) {current $1@D_{16}$};
  
%   \node[anchor=west,text=rewriteOrange!90!black,font=\small] at (.80,.45)
%     {$2w=20$ fixed reductions};
%   \node[anchor=west,text=rewriteGreen!70!black] at (10.50,.45)
%     {expiry: $+0$ reductions};
%   \node[anchor=west,text=rewriteGray!90!black] at (17.00,.45)
%     {tie: output reread};
  
%   \end{tikzpicture}
%   \caption{Chronological adjacent-window reduction.  A generic full DNA4 call
%   processes $L$ candidates using $(k-1)+L$ base ranks and a $(w-1)+L$
%   candidate row; only its final call can be partially active.  The fixed
%   $k=15,w=10$ path constructs eight candidates and reduces them as two
%   successive four-window groups.  Eight is an implementation choice, not a
%   function of $k$ or $w$.  Each group always executes ten shifted reductions,
%   hence twenty per full fast call.  Expiry adds no minimum rescan; tie handling
%   only rereads records for output.}
%   \label{fig:rewrite-chronological-call}
%   \end{figure}
  
%   \begin{algorithm}[H]
%   \caption{Generic chronological adjacent-window kernel}
%   \label{alg:rewrite-generic-chronological}
%   \small
%   \begin{enumerate}
%     \item Let $L$ be the hardware u64 SIMD lane count.  Retain the last $k-1$
%           base ranks, the last $w-1$ candidate records, the current selected
%           record, and a pending-rank count.  The pending array is storage; it is
%           not an algorithmic batch-size parameter.
%     \item Warm up until $k-1$ ranks are retained.  On the DNA4 path, queue every
%           later rank.  Do not invoke the ordinary kernel until exactly $L$
%           ranks are pending.
%     \item For one ordinary invocation, process those $L$ ranks as follows.
%           \begin{enumerate}
%             \item Concatenate the $k-1$ retained ranks and the $L$ pending ranks.
%                   Perform the fixed loop of $k$ shifted vector updates to form
%                   $L$ exact forward and reverse-complement encodings.
%                   Canonicalize them, reject exact palindromes, calculate the
%                   minimap integer hashes, and attach position and strand.
%             \item Prepend the $w-1$ retained candidate records:
%                   \[
%                     D=(x_{a-w+1},\ldots,x_{a-1},
%                        x_a,\ldots,x_{a+L-1}).
%                   \]
%                   For each $r=0,\ldots,w-1$, load
%                   \[
%                     V_r=(D_r,D_{r+1},\ldots,D_{r+L-1})
%                   \]
%                   and update the $L$ running minima with $\leq$.  The result in
%                   lane $\ell$ is the rightmost minimum of the complete adjacent
%                   window ending at $x_{a+\ell}$.
%             \item Align the saved selections as
%                   $(\mathit{current},m_0,\ldots,m_{L-2})$.  Emit the preceding
%                   selection when the entering candidate wins or when that
%                   selection expires.  Expiry uses $m_\ell$, which was already
%                   calculated by the fixed $w$-step reduction.
%             \item If the first-window or expiry gate admits a possible
%                   equal-hash case, reread the relevant window records.  A
%                   confirmed raw-distinct tie uses the chronological event
%                   matrix and compressed output path.  This reread constructs
%                   output records; it never recalculates a window minimum.
%             \item Advance by exactly $L$:
%                   \[
%                     \mathit{candidate\ tail}\gets
%                       (D_L,\ldots,D_{L+w-2}),\qquad
%                     \mathit{current}\gets m_{L-1}.
%                   \]
%                   Retain the last $k-1$ ranks and clear the pending count.
%           \end{enumerate}
%     \item On the ambiguity-capable DNA5 path, process any pending valid ranks
%           before an ambiguity, clear the exact $k$-mer context, and advance
%           invalid sentinel candidates until the context has been rebuilt.
%     \item At finalization, set $q$ to the positive pending remainder.  If the
%           $k$-mer context is complete, then $q<L$.  Invoke the same routine once
%           with lanes $0,\ldots,q-1$ active and the remaining lanes set to
%           sentinels, then advance the two tails and current selection by $q$.
%           Finally append the saved current record if it is valid.
%   \end{enumerate}
%   \end{algorithm}
  
%   \paragraph{Is the specialized kernel different from the generic kernel?}
%   Yes.  It implements the same chronological adjacent-window algorithm and the
%   same rightmost-minimum and output rules, but replaces the generic data path by
%   a fixed $k=15,w=10$ microkernel.
  
%   \begin{enumerate}
%     \item \textbf{Scheduling.}  One ordinary generic call processes exactly
%           $L$ candidates.  One ordinary fast call stages exactly eight
%           candidates and processes them sequentially as two four-lane groups.
%           The numeral eight is fixed by this implementation; it is not computed
%           from $k$ or $w$.  A final remainder uses one masked generic group or
%           one or two partially active four-lane fast groups.
%     \item \textbf{$k$-mer construction.}  The generic path performs $k$ shifted
%           u64 vector updates.  The specialization constructs eight K15
%           encodings with fixed SWAR operations from 22 meaningful ranks in a
%           zero-initialized 24-byte scratch array.  The last two bytes are load
%           padding, not candidates.
%     \item \textbf{Window reduction.}  The generic path uses one
%           $(w-1)+L$ record row and exactly $w$ shifted reductions.  Each fast
%           group fixes this to $9+4=13$ records and ten reductions.  Therefore a
%           full fast call performs $10+10=20$ reductions.
%     \item \textbf{Vector representation.}  The generic path carries hash, raw
%           record, and ordinal in u64 vectors.  The specialization carries four
%           u32 hashes but splits the four raw records and ordinals across pairs
%           of u64x2 vectors.
%     \item \textbf{Domain and fallback.}  The specialization requires DNA4,
%           packed-u32 input, the admitted minimizer layout, and a u32x4/u64x2
%           non-scalar target.  Otherwise the selector uses the generic u64
%           kernel.  Only the generic DNA5 route handles ambiguous symbols by
%           advancing invalid candidates while rebuilding the $k$-mer context.
%   \end{enumerate}
  
%   \subsection{Algorithmic differences: fixed and probability-dependent work}
  
%   Let $N$ be the number of useful candidate positions and let $L$ be the one
%   lane-count symbol used for all three methods.  Thus
%   $\lceil N/L\rceil$ is the number of SIMD groups.  A1 and A2 place separated
%   sequence streams in those lanes, whereas A3 places adjacent windows in them.
  
%   Use the requested two-event model.  Let $p$ be the per-lane probability that
%   the saved minimum expires, and let $r>p$ be the per-lane probability that a
%   window contains a repeated minimum hash.  To obtain a closed form from only
%   $p$ and $r$, assume these events are independent within a lane and assume
%   independent lanes.  This is a deliberately simplified estimate because
%   adjacent A3 lanes share $w-1$ candidates.  Define
%   \[
%     \pi_p=1-(1-p)^L,
%     \qquad
%     \pi_{pr}=1-(1-pr)^L .
%   \]
%   The first quantity estimates the probability that some lane expires.  The
%   second estimates the probability that some lane both expires and contains a
%   repeated minimum hash.  A raw-distinct output tie is a subset of the latter
%   event, so terms using $\pi_{pr}$ below are upper bounds.  Notice that $r>p$
%   does not make the tie branch more frequent than expiry: $pr\leq p$, hence
%   $\pi_{pr}\leq\pi_p$.  One-time startup and final-tail effects are included in
%   the additive $O(k+w)$ term.
  
%   \begin{table}[H]
%   \centering
%   \scriptsize
%   \begin{tabularx}{\linewidth}{p{.11\linewidth}p{.25\linewidth}
%                                p{.29\linewidth}X}
%   \toprule
%   Method & Fixed-schedule work & Probability-dependent work
%          & Expected and worst-case work \\
%   \midrule
%   A1: naive SIMD rescan &
%   One constant-cost advance per group after $O(k+w)$ setup. &
%   With probability $\pi_p$, the shared slow path performs one $w$-slot minimum
%   rescan and one $w$-slot equal-minimum scan.  The value of $r$ changes output,
%   not this schedule. &
%   \[
%    \begin{aligned}
%    \mathbb{E}[T_{\mathrm{A1}}]=O\biggl(&k+w\\[-2pt]
%      &+\left\lceil\frac NL\right\rceil(1+2w\pi_p)\biggr).
%    \end{aligned}
%   \]
%   Worst case:
%   $O\!\left(k+w+\lceil N/L\rceil w\right)$. \\
%   \addlinespace
%   A2: rolling ntHash and prefix--suffix minima &
%   Rolling hash update and two aggregate merges per group; one deterministic
%   $(w-1)$-merge suffix rebuild every $w$ groups. &
%   None.  Expiry and ties do not change the instruction schedule. &
%   Minimum maintenance performs
%   \[
%     \begin{aligned}
%     &2\left\lceil\frac NL\right\rceil\\[-2pt]
%     &\quad +(w-1)\left\lfloor
%         \frac{\lceil N/L\rceil}{w}\right\rfloor
%     \end{aligned}
%   \]
%   merges.  Expected and worst case:
%   $O(k+w+N/L)$. \\
%   \addlinespace
%   A3: generic chronological reduction &
%   Every group performs $k$ shifted exact-$k$-mer construction steps and $w$
%   shifted adjacent-window reductions:
%   \[
%     \begin{aligned}
%     \Theta\biggl(&k+w\\[-2pt]
%       &+\left\lceil\frac NL\right\rceil(k+w)\biggr).
%     \end{aligned}
%   \]
%   Leading term: $\Theta(N(k+w)/L)$. &
%   An expiry-gated group adds a $w$-slot tie-verification pass.  A confirmed tie
%   adds an event construction, transpose, ordering, and compression path costing
%   $O(L(w+L))$ per affected group.  Neither path recomputes a minimum. &
%   \[
%    \begin{aligned}
%    \mathbb{E}[T_{\mathrm{A3}}]=O\biggl(
%      &k+w\\[-2pt]
%      &+\frac NL[k+w(1+\pi_p)]\\[-2pt]
%      &+N(w+L)\pi_{pr}\biggr).
%    \end{aligned}
%   \]
%   Worst case:
%   $O\!\left(k+N[k/L+w+L]\right)$. \\
%   \bottomrule
%   \end{tabularx}
%   \caption{Fixed and conditional source-level work using only $N,L,w,k,p,r$.
%   The table counts vector-loop and tie-layout work, not machine instructions.}
%   \label{tab:rewrite-fixed-variable-work}
%   \end{table}
  
%   For the fixed K15/W10 specialization, substitute $k=15$, $w=10$, and $L=4$.
%   A full eight-candidate call is simply two consecutive four-lane groups, so
%   the sustained stream still contains $\lceil N/L\rceil$ groups.  The numeral
%   eight never enters the asymptotic expression.  Under the two-probability
%   model above, the following is a conservative upper estimate:
%   \[
%     O\!\left(
%       k+w+\left\lceil\frac{N}{2L}\right\rceil
%       +\left\lceil\frac NL\right\rceil
%       \left[
%         w(1+\pi_p+2\pi_{pr})+L(w+L)\pi_{pr}
%       \right]
%     \right).
%   \]
%   The first ceiling counts one fixed packed K15 construction per pair of
%   four-lane groups.  The bracket contains the always-run minimum reduction, the
%   expiry hash prefilter, the tie raw-verification and event-enumeration passes,
%   and the rare event layout.  With $k=15,w=10,L=4$ all of these are constants,
%   so expected and worst-case time are both $\Theta(N)$.
  
%   For one full fast call, the minimum loop itself always performs
%   \[
%     \underbrace{10}_{\text{first group}}
%     +\underbrace{10}_{\text{second group}}
%     =20
%   \]
%   shifted hash comparisons.  Its expected count of ten-slot reduction or
%   verification loops is bounded by
%   \[
%     20(1+\pi_p+2\pi_{pr}),
%     \qquad
%     \pi_p=1-(1-p)^4,
%     \qquad
%     \pi_{pr}=1-(1-pr)^4.
%   \]
%   The two $\pi_{pr}$ terms conservatively charge both raw verification and
%   event enumeration whenever the repeated-hash gate succeeds; the source can
%   reject the event before the latter pass.  The fixed twenty reduction steps
%   load 20 u32 hash vectors, 40 u64 raw-record
%   vectors, and 40 u64 ordinal vectors.  Expiry adds no minimum-reduction load;
%   it can only trigger the separate verification work shown above.
  
%   Algorithmically, A2 has the best guaranteed scaling in $k$ and $w$: after
%   initialization it is $O(N/L)$ and its schedule is independent of $p$ and $r$.
%   A1 can execute fewer minimum operations when expiry is extremely rare, but
%   its worst case grows as $Nw/L$.  Generic A3 deliberately pays
%   $N(k+w)/L$ fixed work to calculate complete adjacent windows.  These orders do
%   not predict fixed-$k$, fixed-$w$ wall-clock speed, where the specialized
%   packing and the different metadata operations can dominate.
  

% \begin{figure}[H]
%   \centering
%   \begin{tikzpicture}[
%     x=.62cm,
%     y=.58cm,
%     font=\footnotesize,
%     >=Latex,
%     work cell/.style={draw=black!75,minimum width=.47cm,
%       minimum height=.46cm,inner sep=0pt},
%     retained/.style={work cell,fill=rewriteGray!20},
%     entering/.style={work cell,draw=rewriteBlue!90!black,
%       fill=rewriteBlue!14},
%     entering two/.style={work cell,draw=rewriteBlue!90!black,
%       densely dashed,fill=rewriteBlue!7},
%     padding/.style={work cell,draw=rewriteGray!80,dashed,
%       fill=rewriteGray!7,text=rewriteGray!85!black},
%     picked/.style={work cell,draw=rewriteOrange!90!black,
%       very thick,fill=rewriteOrange!20},
%     passed tie/.style={work cell,draw=rewriteOrange!90!black,
%       densely dashed,fill=rewriteOrange!7},
%     result/.style={draw=rewriteOrange!90!black,very thick,
%       fill=rewriteOrange!15,rounded corners,align=center,
%       minimum height=.62cm,inner sep=3pt},
%     output/.style={draw=rewriteGreen!80!black,very thick,
%       fill=rewriteGreen!13,rounded corners,align=center,
%       minimum height=.68cm,inner sep=4pt},
%     box/.style={draw=black!70,rounded corners,fill=black!2,
%       align=center,inner sep=4pt},
%     blue box/.style={box,draw=rewriteBlue!85!black,
%       fill=rewriteBlue!10},
%     orange box/.style={box,draw=rewriteOrange!90!black,
%       fill=rewriteOrange!12},
%     green box/.style={box,draw=rewriteGreen!80!black,
%       fill=rewriteGreen!11},
%     note/.style={draw=black!60,rounded corners,align=left,
%       inner sep=4pt,font=\scriptsize},
%     flow/.style={-{Latex[length=1.8mm]},thick}
%   ]
%     % ------------------------------------------------------------------------
%     % (a) Parameter anatomy and base staging.
%     % ------------------------------------------------------------------------
%     \node[anchor=west,font=\bfseries] at (0,28.00)
%       {(a) Parameters: batching width is a SIMD choice, not a function of $k$ or $w$};
  
%     \node[note,text width=13.45cm,anchor=center] at (11.15,26.65) {
%       $L$ is the hardware lane count of one reduction group.  $B$ is the nominal
%       candidate capacity of one outer block.  $A$ is the actual active count,
%       $1\leq A\leq B$.  The active block uses $G=\lceil A/L\rceil$ groups.
%     };
  
%     \node[blue box,text width=6.55cm,font=\scriptsize] at (5.45,23.35) {
%       \textbf{Generic chronological kernel}\\[-1pt]
%       $L=L_{64}\in\{2,4,8,16\}$,\quad $B=L$,\quad $G=1$\\
%       $K_{\max}=31$,\quad target-specific $L_{\max}\leq16$\\
%       base allocation: $(K_{\max}-1)+L_{\max}\leq46$ bytes\\
%       base footprint: $(k-1)+L$ loaded, $(k-1)+A$ meaningful\\
%       three candidate arrays: $(W_{\mathrm{cap}}-1)+L_{\max}$ u64 each\\
%       reduction row: $(w-1)+L$, with $A$ live lanes\\
%       $W_{\mathrm{cap}}=w$ when static or $256$ when runtime
%     };
%     \node[orange box,text width=5.80cm,font=\scriptsize] at (16.45,23.35) {
%       \textbf{K15/W10/B8 specialization}\\[-1pt]
%       $B=8$,\quad $L=L_{32}=4$,\quad full block: $G=B/L=2$\\
%       full-block meaningful ranks: $14+8=22$\\
%       candidate row: $9+4=13$ records\\
%       minimum work: $10$ shifted steps/group
%     };
  
%     \node[anchor=east] at (2.05,20.35) {generic staging};
%     \node[retained,minimum width=3.65cm] at (5.30,20.35)
%       {$k-1$ retained ranks};
%     \node[entering,minimum width=1.90cm] at (10.25,20.35)
%       {$A$ new};
%     \node[padding,minimum width=1.80cm] at (13.55,20.35)
%       {$L-A$ unused};
  
%     \node[anchor=east] at (2.05,18.85) {fast base scratch};
%     \foreach \i in {0,...,13}
%       \node[retained,minimum width=.34cm,minimum height=.42cm]
%         at (2.38+.61*\i,18.85) {\scriptsize o};
%     \foreach \i in {14,...,21}
%       \node[entering,minimum width=.34cm,minimum height=.42cm]
%         at (2.38+.61*\i,18.85) {\scriptsize n};
%     \foreach \i in {22,23}
%       \node[padding,minimum width=.34cm,minimum height=.42cm]
%         at (2.38+.61*\i,18.85) {\scriptsize 0};
%     \node[anchor=south] at (6.35,19.27) {$k-1=14$ retained};
%     \node[anchor=south,text=rewriteBlue!90!black] at (12.95,19.27)
%       {$B=8$ entering};
%     \node[anchor=south,text=rewriteGray!85!black] at (16.05,19.27)
%       {2 padding};
%     \node[note,text width=3.15cm,anchor=west] at (17.15,18.85) {
%       A full block uses $22$ meaningful ranks and two padding bytes.  The equality
%       $24=k+w-1$ is accidental: the biological window span and this
%       24-byte safety-padded allocation are different quantities.  A partial
%       block has only $14+A$ meaningful ranks in the same zeroed array.
%     };
  
%     % ------------------------------------------------------------------------
%     % (b) One B8 block, split into two sequential four-lane groups.
%     % ------------------------------------------------------------------------
%     \node[anchor=west,font=\bfseries] at (0,17.45)
%       {(b) One outer block: $B=8$ new candidates become two sequential $L=4$ groups};
  
%     \node[anchor=east] at (2.15,16.25) {hash records};
%     \foreach \i/\v in {0/8,1/6,2/5,3/9,4/4,5/7,6/4,7/8,8/3}
%       \node[retained] at (2.60+.90*\i,16.25) {\v};
%     \foreach \i/\v in {9/6,10/3,11/5,12/2}
%       \node[entering] at (2.60+.90*\i,16.25) {\v};
%     \foreach \i/\v in {13/7,14/1,15/9,16/1}
%       \node[entering two] at (2.60+.90*\i,16.25) {\v};
%     \foreach \i in {0,...,16}
%       \node[font=\tiny,anchor=north] at (2.60+.90*\i,15.80) {$D_{\i}$};
  
%     \draw[draw=rewriteOrange!90!black,very thick]
%       (2.18,15.35) -- (13.82,15.35);
%     \node[anchor=north,text=rewriteOrange!90!black] at (8.00,15.29)
%       {group 0 row: $D_0,\ldots,D_{12}$, length $(w-1)+L=13$};
%     \draw[draw=rewriteBlue!90!black,very thick,densely dashed]
%       (5.78,14.57) -- (17.42,14.57);
%     \node[anchor=north,text=rewriteBlue!90!black] at (11.60,14.51)
%       {after advancing four: group 1 row $D_4,\ldots,D_{16}$, again 13 records};
  
%     \node[note,text width=3.05cm,anchor=west] at (18.05,15.70) {
%       Here $a$ is the ordinal of the first entering candidate and
%       $D_i=x_{a-w+1+i}$, so $D_9=x_a$.  The notation $v@D_i$ means hash value
%       $v$ at record $D_i$.  Gray is the saved tail and both blue styles are
%       entering groups.  State advances between them.
%     };
  
%     % ------------------------------------------------------------------------
%     % (c) The ten shifted comparisons in group 0.
%     % ------------------------------------------------------------------------
%     \node[anchor=west,font=\bfseries] at (0,13.35)
%       {(c) One group iteration: all $w=10$ shifted SIMD comparisons always run};
%     \node[anchor=east] at (2.15,12.35) {step $r$};
%     \foreach \c in {0,...,9}
%       \node[font=\scriptsize] at (2.80+.90*\c,12.35) {\c};
%     \node[font=\scriptsize,align=center] at (13.45,12.35) {rightmost\\minimum};
  
%     \node[anchor=east,font=\scriptsize,inner sep=1pt] at (2.15,11.35)
%       {$\ell=0$: end $a$};
%     \foreach \c/\v in {0/8,1/6,2/5,3/9,4/4,5/7,6/4,7/8,8/3,9/6}
%       \node[work cell] at (2.80+.90*\c,11.35) {\v};
%     \node[picked] at (2.80+.90*8,11.35) {3};
%     \node[result,minimum width=1.35cm] at (13.45,11.35) {$3@D_8$};
  
%     \node[anchor=east,font=\scriptsize,inner sep=1pt] at (2.15,10.35)
%       {$\ell=1$: end $a+1$};
%     \foreach \c/\v in {0/6,1/5,2/9,3/4,4/7,5/4,6/8,7/3,8/6,9/3}
%       \node[work cell] at (2.80+.90*\c,10.35) {\v};
%     \node[passed tie] at (2.80+.90*7,10.35) {3};
%     \node[picked] at (2.80+.90*9,10.35) {3};
%     \node[result,minimum width=1.35cm] at (13.45,10.35) {$3@D_{10}$};
  
%     \node[anchor=east,font=\scriptsize,inner sep=1pt] at (2.15,9.35)
%       {$\ell=2$: end $a+2$};
%     \foreach \c/\v in {0/5,1/9,2/4,3/7,4/4,5/8,6/3,7/6,8/3,9/5}
%       \node[work cell] at (2.80+.90*\c,9.35) {\v};
%     \node[passed tie] at (2.80+.90*6,9.35) {3};
%     \node[picked] at (2.80+.90*8,9.35) {3};
%     \node[result,minimum width=1.35cm] at (13.45,9.35) {$3@D_{10}$};
  
%     \node[anchor=east,font=\scriptsize,inner sep=1pt] at (2.15,8.35)
%       {$\ell=3$: end $a+3$};
%     \foreach \c/\v in {0/9,1/4,2/7,3/4,4/8,5/3,6/6,7/3,8/5,9/2}
%       \node[work cell] at (2.80+.90*\c,8.35) {\v};
%     \node[picked] at (2.80+.90*9,8.35) {2};
%     \node[result,minimum width=1.35cm] at (13.45,8.35) {$2@D_{12}$};
  
%     \pgfmathsetmacro{\vx}{2.80+.90*3}
%     \draw[draw=rewriteBlue!90!black,very thick,rounded corners]
%       ({\vx-.43},7.91) rectangle ({\vx+.43},11.79);
  
%     \node[orange box,text width=7.65cm,minimum height=.78cm,
%           font=\tiny,inner sep=2pt] at (7.90,7.00) {
%       At every column $r$, compare
%       $V_r[\ell]=h(D_{r+\ell})$ with lane $\ell$'s running minimum.
%       The same $\leq$ mask selects hash, raw record, and ordinal.  The outlined
%       column is $V_3=(D_3,D_4,D_5,D_6)$.
%     };
%     \node[note,text width=4.05cm,anchor=west] at (15.55,9.85) {
%       Each row is one complete adjacent window.  Dashed orange cells are
%       earlier equal minima.  Because the update uses $\leq$, lanes 1 and 2
%       finish at the later $3@D_{10}$.
%     };
  
%     % ------------------------------------------------------------------------
%     % (d) State update, group 1, and all additional work.
%     % ------------------------------------------------------------------------
%     \node[anchor=west,font=\bfseries] at (0,5.85)
%       {(d) Finish the group, advance the state, and distinguish minimum work from output work};
  
%     \node[orange box,text width=2.20cm] (gzero) at (1.80,4.05)
%       {group 0 minima\\$(3@D_8,3@D_{10},$\\$3@D_{10},2@D_{12})$};
%     \node[green box,text width=2.45cm,font=\scriptsize] (events) at (6.40,4.05)
%       {steady lane: emit the preceding selection if entering hash $\leq$ saved
%        hash, or the saved ordinal expires};
%     \node[blue box,text width=2.65cm] (advance) at (11.40,4.05)
%       {$\text{tail}\gets D_4,\ldots,D_{12}$\\
%        $\text{current}\gets2@D_{12}$\\
%        $\text{count}\gets\text{count}+4$};
%     \node[blue box,text width=1.85cm] (gone) at (16.15,4.05)
%       {group 1\\same ten comparison steps};
%     \node[orange box,text width=2.55cm,font=\scriptsize] (gtwo) at (20.55,4.05)
%       {group 1 minima\\$(2@D_{12},1@D_{14},$\\
%        $1@D_{14},1@D_{16})$\\
%        tail $\gets D_8,\ldots,D_{16}$\\current $\gets1@D_{16}$};
%     \draw[flow] (gzero)--(events);
%     \draw[flow] (events)--(advance);
%     \draw[flow] (advance)--(gone);
%     \draw[flow] (gone)--(gtwo);
  
%     \node[orange box,text width=5.30cm,minimum height=.92cm,font=\scriptsize] (fixed)
%       at (4.90,1.00) {
%         \textbf{Always-run minimum work}\\
%         $G\,w=2\cdot10=20$ shifted SIMD reduction steps per full B8 block,
%         hence $80$ lane-wise hash comparisons.  Each step loads one u32x4 hash
%         vector plus two raw and two ordinal u64x2 metadata vectors.
%     };
%     \node[output,text width=4.15cm,minimum height=.92cm] (expiry)
%       at (14.20,1.00) {
%         \textbf{Minimum expires}\\
%         use the minimum already calculated above:
%         \emph{zero additional minimum-reduction steps}
%     };
%     \node[blue box,text width=2.35cm,minimum height=.92cm] (basetail)
%       at (20.25,1.00) {
%         end of block\\
%         base tail $\gets$ last $k-1=14$ ranks
%     };
  
%     \node[note,dashed,text width=8.50cm,anchor=center] (rare)
%       at (11.20,-1.80) {
%         \textbf{Rare tie-output branch:}
%         possible first-window or expiry tie $\longrightarrow$
%         $w$ hash-prefilter slots $\longrightarrow$ duplicate hash
%         $\longrightarrow$ $w$ raw-verification slots $\longrightarrow$
%         raw-distinct tie $\longrightarrow$ $w$ event-enumeration and ordering
%         slots.  No step recomputes the minimum.
%     };
%     \draw[flow,dashed]
%       (events.south) -- ++(0,-.10)
%       -| ([xshift=-2.4cm]rare.west) -- (rare.west);
%   \end{tikzpicture}
%   \caption{One steady-state, full outer-block iteration of the chronological
%   $k=15,w=10,B=8$ fast path, with generic dimensions shown in panel (a).
%   The block first constructs eight adjacent candidate records.  Two sequential
%   $L=4$ groups each concatenate four entering records with the retained
%   $w-1=9$ records.  This is a batched full-window reduction: every group
%   recomputes four complete minima with exactly $w=10$ shifted loads.  Group state
%   advances before the second group, so the full block calculates eight adjacent
%   window minima using twenty fixed reduction steps.  There is no saved-minimum
%   fast path and expiry never starts another minimum rescan.  A possible
%   first-window or expiry tie can initiate a window probe even when the probe
%   disproves the tie.  Confirmed ties add output enumeration, never minimum
%   recalculation.}
%   \label{fig:rewrite-exact-iteration-upgraded}
%   \end{figure}
  
%   \begin{algorithm}[H]
%   \caption{Generic chronological adjacent-window kernel}
%   \label{alg:rewrite-generic-chronological}
%   \small
%   \begin{enumerate}
%     \item Let $L$ be the hardware u64 SIMD lane count and set the nominal
%           full-block capacity to $B=L$.  Let $A\leq L$ be the active count in
%           the present block.  Initialize the $w-1$ candidate tail and saved
%           current selection to maximum sentinels.
%     \item Warm up until $k-1$ base ranks have been retained.  Packed DNA4 fills
%           this rank tail without advancing the candidate state.  Its first real
%           $k$-mer has ordinal zero and final-base position $k-1$.  DNA5 instead
%           advances one invalid candidate for every warm-up rank so valid-run
%           lengths and explicit positions remain aligned.  Neither route forms a
%           real $k$-mer during warm-up.
%     \item Take $A$ new consecutive positions and concatenate their ranks with
%           the $k-1$ base tail in a zero-initialized scratch array.  Shifted SIMD
%           loads touch $(k-1)+L$ entries, while only $(k-1)+A$ ranks are
%           biologically meaningful.
%     \item Form up to $A$ exact forward and reverse-complement $k$-mer encodings
%           from the shifted rank loads.  Canonicalize each active lane, reject
%           exact palindromes, calculate its minimap integer hash, and attach its
%           ordinal, final-base position, and strand.  Replace every inactive
%           candidate hash, raw record, and ordinal by its maximum sentinel.
%     \item Form the logical candidate row
%           \[
%             D=(\underbrace{x_{a-w+1},\ldots,x_{a-1}}_{w-1\text{ retained}},
%                \underbrace{x_a,\ldots,x_{a+L-1}}_{L\text{ SIMD slots}}),
%           \]
%           whose full width is $(w-1)+L$.  Inactive final-block slots contain
%           sentinels.
%     \item Initialize the $L$ running minima to sentinels.  For every
%           $r=0,\ldots,w-1$, unconditionally load
%           \[
%             V_r=(D_r,D_{r+1},\ldots,D_{r+L-1})
%           \]
%           and update every lane with $\leq$.  Use the same comparison mask to
%           select the winning raw record and ordinal.  This produces up to $L$
%           adjacent-window rightmost minima after exactly $w$ shifted steps.
%           Denote the lane results by $m_0,\ldots,m_{L-1}$.
%     \item Align the selection preceding each lane as
%           $(\textit{current},m_0,\ldots,m_{L-2})$.  For each steady full-window
%           lane with a valid predecessor, emit that preceding saved record when
%           the entering hash is less than or equal to its hash.  Also emit it
%           when its ordinal expires while the entering hash is larger.  Expiry
%           consumes the already-calculated lane minimum and does not repeat the
%           $w$-step reduction.  Compress ordinary events in lane order.
%     \item If the first-window or expiry gate indicates a possible tie, reread
%           the relevant window slots to verify it.  If raw-distinct equal-minimum
%           records exist, construct one event row per lane, transpose the event
%           matrix, and remove only sentinels.  The first full window has no
%           ordinary emit-old event.  Only its additional equal-minimum records
%           use this special rule, while a unique minimum remains saved for a
%           later transition or final flush.  No tie pass recomputes a minimum.
%     \item Advance by the actual $A$ positions:
%           \[
%             \textit{candidate tail}\gets(D_A,\ldots,D_{A+w-2}),\qquad
%             \textit{current}\gets m_{A-1},
%           \]
%           and retain the last $k-1$ base ranks from the staging region.  In the
%           generic DNA5 route, an ambiguity first processes pending real
%           candidates, clears the exact $k$-mer context, and advances one invalid
%           candidate for the ambiguous symbol.  Subsequent warm-up ranks also
%           advance invalid candidates until the context is rebuilt.  This does
%           not create an artificial final flush.
%     \item At finalization, process a nonempty partial block only when $k$-mer
%           warm-up is complete.  Then perform the final-current flush, which
%           appends the saved record only if it is valid.
%   \end{enumerate}
%   \end{algorithm}
  
%   \paragraph{Is the specialized kernel different from the generic kernel?}
%   Yes.  Both kernels use adjacent windows, rightmost replacement on $\leq$, the
%   same minimap-style event rules, and no expiry-triggered minimum rescan.  The
%   specialization replaces the generic data path with a fixed K15/W10
%   microkernel:
%   \begin{enumerate}
%     \item \textbf{Batch geometry.}  The generic kernel uses one u64 group with
%           $B=L_{64}$.  The specialization hard-codes $B=8$ and runs two
%           sequential u32x4 hash groups, while raw records and ordinals use
%           u64x2 vectors.
%     \item \textbf{$k$-mer staging.}  The generic kernel incorporates $k$ shifted
%           rank vectors into exact encodings and uses a maximum-capacity scratch
%           allocation.  For a full B8 block, the specialization packs 22
%           meaningful K15/B8 ranks into a fixed 24-byte packed-word staging array
%           with two safety-padding bytes.  A partial block copies only $14+A$
%           meaningful ranks into the same zero-initialized allocation.
%     \item \textbf{Window reduction.}  Both retain $w-1$ records and perform $w$
%           unconditional shifted steps.  The specialization fixes these values
%           to a 13-record group row and ten steps.  A full B8 block has two groups
%           and therefore twenty steps.
%     \item \textbf{Metadata movement and tie output.}  A generic step loads one
%           u64 vector each for hash, raw record, and ordinal.  A specialized step
%           uses one u32x4 hash load plus two u64x2 raw and two u64x2 ordinal loads.
%           Its rare tie path uses fixed four-lane prefilter, verification, event
%           layout, and compressed-store operations.
%     \item \textbf{Admission and fallback.}  The fast path requires DNA4,
%           $k=15,w=10$, packed-u32 input, an admitted output layout, a
%           u32x4/u64x2 non-scalar target, and
%           $\textit{consumed\_bp}\leq\texttt{UINT32\_MAX}$.  When fast-path
%           admission fails, the selector chooses the generic-u64 kind, which
%           still requires an admitted multi-lane u64 target.  The fragmented
%           driver never opts into the fast kind.
%     \item \textbf{Ambiguous symbols.}  The specialization is DNA4-only.  The
%           generic DNA5 route processes pending candidates at an ambiguity,
%           clears the $k$-mer context, and advances invalid candidates without a
%           boundary final flush.
%   \end{enumerate}
  
%   \subsection{Algorithmic differences: fixed and probability-dependent work}
  
%   Let $N$ be the number of useful candidate positions and use one symbol $L$
%   for logical group width throughout.  The symbolic comparison holds this $L$
%   equal across methods.  Let $B$ be the number of consecutive candidates formed
%   by one A3 outer block, so a full block has $G=B/L$ groups for the implemented
%   divisible sizes.  Concrete implementation counts substitute $L=8$ logical
%   separated streams for A2 and $L=4$, $B=8$ for the A3 fast path.
  
%   Assume lanes are mutually independent and, within each lane, assume that
%   expiry and the relevant tie event are independent.  Let $p$ be the per-lane
%   probability that the saved minimum expires and let $r$ be the per-lane
%   probability of a relevant raw-distinct equal-minimum hash tie.  For the
%   illustrative scenario below, assume $r>p$, so
%   $0\leq p<r\leq1$.  This ordering is a model input, not a property guaranteed
%   by the kernels.  For a logical SIMD stage, define
%   \[
%     \pi_p=1-(1-p)^L,
%     \qquad
%     \pi_{pr}=1-(1-pr)^L.
%   \]
%   Here $\pi_p$ is the probability that at least one lane expires.  The second
%   quantity is the probability that at least one lane both expires and has a
%   tie, under the stated independence model.  First-window setup and collector
%   writes are omitted.  A ``stage'' below is an algorithmic vector-loop stage,
%   not necessarily one machine instruction.
  
%   \begin{table}[H]
%   \centering
%   \scriptsize
%   \begin{tabularx}{\linewidth}{p{.12\linewidth}p{.21\linewidth}p{.21\linewidth}X}
%   \toprule
%   Method & Fixed schedule & Probability-dependent schedule & Expected time and worst case \\
%   \midrule
%   A1: naive SIMD rescan &
%   After $O(k)$ warm-up, one rolling $k$-mer-code update, minimap hash evaluation,
%   ring update, and entering-versus-saved comparison advance $L$ separated
%   candidates. &
%   If any lane expires, the code always performs one $w$-slot minimum rescan and
%   a second $w$-slot tie scan, even when that scan emits nothing.  This occurs
%   with probability $\pi_p$.  The tie probability $r$ changes output count, not
%   the number of scans. &
%   \[
%     \begin{aligned}
%     &\mathbb{E}[T_{\mathrm{A1}}]
%       =O\!\biggl(k+\frac NL\\[-2pt]
%     &\qquad{}\cdot(1+2w\pi_p)\biggr).
%     \end{aligned}
%   \]
%   Worst case: $O(k+Nw/L)$. \\
%   \addlinespace
%   A2: rolling ntHash with prefix--suffix minima &
%   After $O(k)$ initialization, every $L$-candidate advance performs rolling
%   hash work and two aggregate-minimum merges.  Every $w$ advances add one fixed
%   $(w-1)$-merge suffix rebuild. &
%   Neither expiry nor a hash tie starts a window reread.  Ties update the two
%   stored endpoints and may change output count, but not the minimum schedule. &
%   With $T=N/L$ logical advances, the minimum structure executes
%   \[
%     2T+(w-1)\left\lfloor\frac Tw\right\rfloor
%   \]
%   merges.  Total and worst case: $O(k+N/L)$. \\
%   \addlinespace
%   A3: chronological adjacent-window reduction &
%   The generic path performs $O(k)$ exact-code stages per $B$-candidate block,
%   then exactly $w$ shifted reductions per $L$-window group.  Its fixed work is
%   $O(Nk/B+Nw/L)$. &
%   An expiry-gated group adds one $w$-slot verification pass.  If that probe also
%   contains a raw-distinct equal minimum, the padded event construction,
%   transpose, ordering, and compression cost $O(L(w+L))$ per group when $L$ is
%   an asymptotic variable.  The fast path first adds a $w$-slot raw-verification
%   pass. &
%   \[
%     \begin{aligned}
%     &\mathbb{E}[T_{\mathrm{A3,g}}]\\[-2pt]
%     &\quad=O\!\biggl(\frac{Nk}{B}
%       +\frac{Nw}{L}(1+\pi_p)\\[-2pt]
%     &\qquad{}+N(w+L)\pi_{pr}\biggr),
%     \end{aligned}
%   \]
%   \[
%     \begin{aligned}
%     &\mathbb{E}[T_{\mathrm{A3,f}}]\\[-2pt]
%     &\quad=O\!\biggl(\frac{N}{B}
%       +\frac{Nw}{L}(1+\pi_p)\\[-2pt]
%     &\qquad{}+N(w+L)\pi_{pr}\biggr).
%     \end{aligned}
%   \]
%   Worst case: $O(Nk/B+Nw/L+N(w+L))$ generically and
%   $O(N/B+Nw/L+N(w+L))$ for fixed K15. \\
%   \bottomrule
%   \end{tabularx}
%   \caption{Source-level work for A1, A2, and A3 using one lane-count symbol
%   $L$.  A1 and A2 lanes are separated streams, whereas an A3 group contains
%   adjacent windows.  Probability changes conditional work but not the stated
%   worst-case orders.}
%   \label{tab:rewrite-fixed-variable-work}
%   \end{table}
  
%   For the concrete A3 specialization, $k=15$, $w=10$, $B=8$, $L=4$, and
%   $G=2$.  Its fixed minimum work is
%   \[
%     Gw=20
%   \]
%   shifted reductions per block.  Excluding matrix ordering and compression, its
%   exact expected count of ten-slot reduction or scan loops is
%   \[
%     Gw(1+\pi_p+2\pi_{pr})
%     =20(1+\pi_p+2\pi_{pr}),
%     \qquad
%     \pi_p=1-(1-p)^4,
%     \qquad
%     \pi_{pr}=1-(1-pr)^4.
%   \]
%   Confirmed ties additionally incur
%   $O(G L(w+L)\pi_{pr})$ expected event-matrix ordering and compression work in
%   the $L$-variable model.  With the implemented fixed $L=4$, this is
%   $O(G(w+L)\pi_{pr})$.
%   The fixed twenty reductions perform 20 u32 hash loads and comparisons, 40 u64
%   raw-record loads, and 40 u64 ordinal loads.  A2 with $L=8$ and $w=10$
%   performs $3w-1=29$ logical aggregate merges per ten advances, covering 80
%   candidates.  These counts are not instruction-equivalent because an A2
%   aggregate merge is heavier than one A3 hash comparison.
  
%   The corresponding asymptotic working memories are
%   \[
%     M_{\mathrm{A1}}=O(wL),
%     \qquad
%     M_{\mathrm{A2}}=O((k+w)L),
%     \qquad
%     M_{\mathrm{A3}}=O(k+B+wL+L^2).
%   \]
  
%   A2 has the lowest guaranteed asymptotic minimum-maintenance cost as $w$
%   grows: it is amortized $O(1)$ per logical candidate and its schedule is
%   independent of $p$ and $r$.  A1 has expected cost
%   $O(1+2w\pi_p)$ per $L$-candidate advance and an $O(w)$ slow-path burst.  A3
%   pays $O(w/L)$ fixed reductions per candidate,
%   $O(w\pi_p/L)$ expected expiry-verification work,
%   $O((w+L)\pi_{pr})$ expected tie ordering, and the generic $O(k/B)$
%   hash-construction term.  Thus A2 has the best guaranteed scaling in $w$ under
%   these assumptions.  This does not rank
%   fixed-$k$, fixed-$w$ wall-clock performance because the kernels move different
%   metadata and implement different score and tie semantics.
  




% \subsection{What \texttt{bounded\_nthash\_kernel.cppm} does}
% \label{sec:rewrite-bounded-source}

% The bounded kernel fixes $k=15$, $w=10$, DNA4 input, and eight logical
% streams.  A physical target with four 32-bit lanes holds those streams in two
% vector registers; a target with eight 32-bit lanes holds them in one.  Logical
% stream $q$ begins at a different packed-sequence offset, owns a disjoint range
% of output windows, and reads right-hand context.  Therefore lanes at one SIMD
% step are far apart in the biological sequence.  They are not $W_j,W_{j+1},
% \ldots$.

% Let $r_i\in\{0,1,2,3\}$ be the DNA4 rank and let
% \[
%   T=(\mathtt{95c60474},\mathtt{62a02b4c},
%      \mathtt{82572324},\mathtt{4be24456})
% \]
% be the four 32-bit table words written in hexadecimal.  With 32-bit rotation
% \(\operatorname{ROL}\), the kernel initializes
% \begin{align}
%   F_i &= \bigoplus_{j=0}^{k-1}
%     \operatorname{ROL}_{7(k-1-j)}\!\left(T[r_{i+j}]\right),\\
%   R_i &= \bigoplus_{j=0}^{k-1}
%     \operatorname{ROL}_{7j}\!\left(T[3-r_{i+j}]\right).
% \end{align}
% If $a=r_i$ leaves and $b=r_{i+k}$ enters, the steady-state update is
% \begin{align}
%   F_{i+1} &= \operatorname{ROL}_{7}(F_i)
%     \oplus \operatorname{ROL}_{7k}(T[a]) \oplus T[b],\\
%   R_{i+1} &= \operatorname{ROR}_{7}(R_i)
%     \oplus \operatorname{ROR}_{7}(T[3-a])
%     \oplus \operatorname{ROL}_{7(k-1)}(T[3-b]).
% \end{align}
% All rotation counts are modulo 32; hence for $k=15$ the two compile-time
% left rotations are $7k\bmod32=9$ and $7(k-1)\bmod32=2$.  The ordering score
% and strand bit are
% \[
%   h_i=(F_i+R_i)\bmod 2^{32},
%   \qquad z_i=\mathbf 1[R_i<F_i].
% \]
% This is a rolling ntHash-like score.  It is not
% $\min(F_i,R_i)$ and it is not minimap's integer permutation.

% For each far-apart stream, the kernel feeds these records into the A2
% prefix--suffix structure.  Its persistent state contains a ring of $k=15$
% symbol vectors, a prefix-minimum aggregate, and a ring of $w=10$
% suffix-minimum aggregates.  Every tenth candidate it rebuilds all suffixes in
% a fixed descending loop; every window is then obtained by combining the
% current prefix with one saved suffix.  Equal scores are merged into an
% aggregate that remembers only the leftmost and rightmost attaining positions.
% The left endpoint is considered for output first and the right endpoint is
% also considered when distinct.  Output suppression compares the complete
% identity $(h,p,z)$, not merely $h$, and only against the immediately preceding
% identity in that logical stream.  Consequently, a persistent tied pair written
% as left--right may be written as left--right again for the next window;
% interior occurrences of a three-way or larger tie are never retained.  The
% first event of streams 1 through 7 is temporarily retained so that the
% duplicated chunk boundary can be suppressed when the eight streams are joined.
% Physical collector order across the lockstep streams is intentionally
% unspecified.  There is no matrix transpose, no conditional minimum rescan, and
% no minimap-style final-current flush in this kernel: every owned full window is
% considered online, after which finalization only reconciles deferred stream
% boundaries.

% The quantity $k+w-1=24$ here is the number of biological bases required to
% finish a minimizer window that starts at a chunk boundary.  Because a packed
% 32-bit word contains 16 DNA4 bases, the iterator rounds that context to two
% words, i.e. 32 decoded bases.  It is not a 24-element hash buffer.  The
% algorithmic state is the separate 15-entry symbol ring and 10-entry minimum
% ring described above.

% The bounded kernel contains no explicit palindrome mask.  This is harmless
% for its only admitted $k$: an exact reverse-complement palindrome over DNA4
% cannot have odd length, because its middle base would have to equal its own
% complement.  Equality $F_i=R_i$ can still occur as a 32-bit hash collision;
% the code retains that candidate and assigns strand zero.

% \subsection{The chronological adjacent-window kernel}
% \label{sec:rewrite-adjacent-correct}

% The following description is specifically of
% \texttt{two\_stack\_nthash}'s chronological DNA4 $k=15,w=10$ fast path in
% \texttt{exact\_k15\_w10.cppm}.  It must not be attributed to
% \texttt{bounded\_nthash\_kernel.cppm} or to the resident-vector route.

% Let $\rho(A)=0$, $\rho(C)=1$, $\rho(G)=2$, and $\rho(T)=3$.  For the $k$-mer
% beginning at $i$, this path constructs the exact two-bit encodings
% \begin{align}
%   F_i &= \sum_{j=0}^{k-1}\rho(s_{i+j})\,2^{2(k-1-j)},\\
%   R_i &= \sum_{j=0}^{k-1}\bigl(3-\rho(s_{i+j})\bigr)\,2^{2j},\\
%   C_i &= \min(F_i,R_i),
%   &z_i&=\mathbf 1[R_i<F_i],
%   &q_i&=\mathbf 1[F_i=R_i].
% \end{align}
% Thus $q_i$ is the explicit reverse-complement-palindrome test that was absent
% from the previous description.  A palindromic lane is changed to an invalid
% maximum sentinel before minimum selection.  Again $q_i$ is necessarily zero
% for valid DNA4 when $k=15$, but the check is present in the executed code and
% the generic route also supports other $k$ values.

% For $k=15$, put $M=2^{30}-1$ and $u_0=C_i$.  The actual minimap integer
% permutation $\phi$ is the following sequence of unsigned 32-bit operations:
% \begin{align}
%   u_1&=(\operatorname{not}u_0+(u_0\ll21))\mathbin{\&}M,
%   &u_2&=u_1\oplus(u_1\gg24),\\
%   u_3&=(u_2+(u_2\ll3)+(u_2\ll8))\mathbin{\&}M,
%   &u_4&=u_3\oplus(u_3\gg14),\\
%   u_5&=(u_4+(u_4\ll2)+(u_4\ll4))\mathbin{\&}M,
%   &u_6&=u_5\oplus(u_5\gg28),\\
%   h_i=\phi(C_i)&=(u_6+(u_6\ll31))\mathbin{\&}M.
% \end{align}
% Let
% \[
%   H_\infty=
%   \operatorname{uint32}\!\left(
%     \texttt{maximum\_record.kmer()}\right)
% \]
% be the hash sentinel supplied by the selected output representation.  A lane
% is usable only when $q_i=0$ and $h_i\ne H_\infty$; otherwise its hash and raw
% record are replaced by maximum sentinels.  This distinction matters: for the
% 30-bit target layout $H_\infty=M$, whereas for the 35-bit query layout its
% 32-bit cast is $2^{32}-1$, so the genuine 30-bit value $h_i=M$ remains valid.
% The stored coordinate is the final base of the $k$-mer, $p_i=i+k-1$, and the
% raw identity includes $(h_i,p_i,z_i)$.  No rolling-hash recurrence is used: up
% to eight overlapping exact encodings are formed together by SWAR packing and
% SIMD bit operations.

% \subsubsection{What ``block'', ``group'', and ``lane'' mean}

% The specialized routine accepts $A\in\{1,\ldots,8\}$ active consecutive
% candidate positions, with a full-block capacity $B=8$.  For a full block, hash
% construction uses a 24-byte array of which only
% \[
%   (k-1)+B=14+8=22
% \]
% bytes contain bases: 14 retained ranks followed by eight new ranks.  The final
% two zero bytes exist only so that three eight-byte SWAR loads remain in bounds.
% For a final partial block, only $14+A$ bytes are meaningful and the remaining
% $10-A$ bytes stay zero in the value-initialized array.  These bytes are
% padding, not sequence context.

% The up to eight candidate hashes are split into at most two \emph{groups} of
% $L=4$, because the admitted fast shape has four 32-bit lanes (and two 64-bit
% lanes for packed output metadata).  Inactive lanes become sentinels.  Before
% reducing one group, the code concatenates the
% $w-1=9$ retained candidate records with its four entering candidates.  The
% minimum buffer therefore contains
% \[
%   (w-1)+L=9+4=13
% \]
% candidate records.  If $a$ is the ordinal of the first entering candidate in
% the group, write this combined row as
% \[
%   D_r=x_{a-w+1+r},\qquad 0\le r<w+L-1.
% \]
% During startup, define $x_j=\bot$ for $j<0$ with
% $h(\bot)=H_\infty$.  A lane represents a complete biological window, and is
% eligible for emission, only when $a+\ell\ge w-1$; earlier lanes are padded
% warm-up reductions.

% For each slot $r=0,\ldots,w-1$, the unaligned hash-vector load produces
% \[
%   V_r[\ell]=h(D_{r+\ell})
%             =h_{a-w+1+r+\ell},
%   \qquad 0\le\ell<L.
% \]
% The ten vectors are reduced lane-wise.  Hence lane $\ell$ sees exactly
% \[
%   \{V_0[\ell],\ldots,V_{w-1}[\ell]\}
%   =\{h_{a-w+1+\ell},\ldots,h_{a+\ell}\},
% \]
% For an active lane after warm-up, this is the complete window ending at
% candidate $a+\ell$.  A group therefore performs four adjacent lane reductions
% representing up to four active complete windows; inactive or warm-up lanes are
% sentinel-padded and are not emitted.  After the first group updates the
% nine-record tail, the second group processes up to four further windows in the
% same way.  This is unrelated to the far-apart lanes in A1, A2, or the resident
% route.

% \begin{figure}[H]
% \centering
% \begin{tikzpicture}[font=\footnotesize,x=1cm,y=1cm]
%   \node[anchor=west,font=\bfseries] at (0,7.55)
%     {(a) Full-block hash construction: 22 meaningful bytes in a 24-byte array};
%   \node[anchor=east] at (1.55,6.85) {base ranks};
%   \foreach \i in {0,...,13}
%     \node[rewrite old,minimum width=.43cm,minimum height=.44cm]
%       at (1.82+.46*\i,6.85) {o};
%   \foreach \i in {14,...,21}
%     \node[rewrite new,minimum width=.43cm,minimum height=.44cm]
%       at (1.82+.46*\i,6.85) {n};
%   \foreach \i in {22,23}
%     \node[rewrite cell,draw=rewriteGray!80,dashed,fill=rewriteGray!8,
%           minimum width=.43cm,minimum height=.44cm]
%       at (1.82+.46*\i,6.85) {0};
%   \draw[decorate,decoration={brace,mirror,amplitude=3pt}]
%     (1.59,6.50)--(7.99,6.50)
%     node[midway,below=4pt] {$k-1=14$ retained bases};
%   \draw[decorate,decoration={brace,mirror,amplitude=3pt}]
%     (8.03,6.50)--(11.71,6.50)
%     node[midway,below=4pt] {$B=8$ new bases};
%   \draw[decorate,decoration={brace,mirror,amplitude=3pt}]
%     (11.75,6.50)--(12.63,6.50)
%     node[midway,below=4pt] {padding};

%   \node[anchor=west,font=\bfseries] at (0,5.35)
%     {(b) One group: 9 retained candidates plus 4 new candidates};
%   \node[anchor=east] at (1.55,4.72) {candidate row};
%   \foreach \i/\v in {0/8,1/6,2/5,3/9,4/4,5/7,6/4,7/8,8/3}
%     \node[rewrite old,minimum width=.52cm] at (1.85+.57*\i,4.72) {\v};
%   \foreach \i/\v in {9/6,10/3,11/5,12/2}
%     \node[rewrite new,minimum width=.52cm] at (1.85+.57*\i,4.72) {\v};
%   \foreach \i in {0,...,12}
%     \node[font=\scriptsize,anchor=north] at (1.85+.57*\i,4.39) {$D_{\i}$};

%   \foreach \ell/\first/\winner/\result in {0/0/8/3,1/1/10/3,2/2/10/3,3/3/12/2} {
%     \pgfmathsetmacro{\yy}{3.72-.62*\ell}
%     \node[anchor=east] at (1.55,\yy) {lane $\ell$};
%     \foreach \j in {0,...,9} {
%       \pgfmathtruncatemacro{\idx}{\first+\j}
%       \pgfmathparse{int(\idx==\winner)}
%       \ifnum\pgfmathresult=1
%         \node[rewrite min,minimum width=.52cm]
%           at (1.85+.57*\idx,\yy) {\result};
%       \else
%         \node[rewrite cell,minimum width=.52cm]
%           at (1.85+.57*\idx,\yy) {};
%       \fi
%     }
%   }
%   \node[rewrite note,text width=3.65cm,anchor=west] at (9.35,2.78)
%     {Hash step $r$ supplies the diagonal
%      $(D_r,D_{r+1},D_{r+2},D_{r+3})$ to lanes 0--3.
%      Lanes 1 and 2 tie at $D_8,D_{10}$; the $\leq$ update selects
%      $D_{10}$ in both.};
% \end{tikzpicture}
% \caption{The two different working sets in the $k=15,w=10,B=8$ fast path.
% Panel (a) is the base scratch used to construct eight exact $k$-mers.  Panel
% (b) is the 13-record candidate buffer used to reduce one four-window group.
% A minimizer window itself spans $k+w-1=24$ biological bases, but neither
% working set is a 24-candidate minimizer buffer.}
% \label{fig:rewrite-exact-buffer-and-group}
% \end{figure}

% The reduction initializes each lane to the maximum sentinel and visits slots
% from left to right using $\leq$.  Therefore a later occurrence replaces an
% earlier occurrence with the same hash:
% \[
%   m_{\ell}^{(-1)}=\bot,\qquad
%   m_{\ell}^{(r)}=
%   \begin{cases}
%     D_{r+\ell}, & h(D_{r+\ell})\le h(m_{\ell}^{(r-1)}),\\
%     m_{\ell}^{(r-1)}, & \text{otherwise},
%   \end{cases}
%   \quad r=0,\ldots,w-1.
% \]
% Thus $m_{\ell}^{(w-1)}$ is the rightmost argmin of the corresponding window.
% The loop over all $w=10$ slots runs for every group.  Each slot performs one
% shifted hash-vector comparison step and loads the corresponding metadata as
% two u64 raw-record vectors and two u64 ordinal vectors.  It is best understood
% as an \emph{unconditional batched full-window reduction}: four fixed lane
% reductions are performed together by ten comparison steps, with
% sentinel padding for inactive or warm-up lanes.  It is not A1's conditional
% rescan, because expiry does not decide whether the reduction runs, and it is
% not A2's prefix--suffix method, because no two saved minima are combined to
% obtain these windows.

% \begin{algorithm}[H]
% \caption{Chronological $k=15,w=10,B=8$ adjacent-window kernel}
% \label{alg:rewrite-exact-adjacent}
% \small
% \begin{enumerate}
%   \item Retain the last 14 DNA4 ranks, the last nine complete candidate
%         records, and the current minimap-style selected record.
%   \item Append $A\leq8$ new ranks to the 14 retained ranks in a zero-initialized
%         24-byte array.  A full block has two padding bytes; a partial block has
%         $10-A$ padding bytes.  The padding makes all three SWAR loads safe.
%   \item Construct up to eight consecutive forward and reverse-complement
%         30-bit encodings.  Canonicalize them, apply the minimap integer
%         permutation, record final-base position and strand, and replace
%         palindromes, inactive lanes, or representation-reserved maximum hashes
%         by sentinels.
%   \item Split the active candidates into at most two chronological groups of
%         four.  For each group, prepend the nine retained candidate records and
%         perform ten shifted hash-vector comparison steps.  Each step also
%         loads two raw-record and two ordinal u64 vectors.  Reduce with $\leq$,
%         so each full-window lane obtains the rightmost minimum.
%   \item In increasing window order, apply the exact minimap event rules
%         described below.  Use two stable compressed stores, one for each u64
%         lane half, for the ordinary event mask; use the event-matrix path only
%         if a raw-distinct equal-hash tie really exists.
%   \item After a group, retain its last nine candidate records.  After the B8
%         block, retain its last 14 base ranks.  At end of input, process a
%         partial block and append the final saved minimum if it is valid.
% \end{enumerate}
% \end{algorithm}

% \subsubsection{Output events, equal hashes, and the matrix path}

% Calculating one rightmost argmin per window is not, by itself, sufficient to
% reproduce minimap's output.  Let $e_t$ be the entering candidate at ordinal
% $t$ and let $c_{t-1}$ be the saved selected record immediately before it.  On
% an ordinary full-window step, the old record becomes an output event when
% \[
%   h(e_t)\le h(c_{t-1})
%   \quad\text{or}\quad
%   \operatorname{ord}(c_{t-1})=t-w.
% \]
% The first condition introduces a new, possibly equal, rightmost selection;
% the second says that the saved identity just expired.  A minimum that persists
% through many adjacent windows is therefore eventually written once, while an
% equal hash at a different position remains a different raw record.  The last
% saved-record flush is called unconditionally at end of input, even if fewer
% than $w$ candidates were seen, and appends a record only when that saved state
% is valid.

% Most groups contain no complicated tie.  Their old records are selected by a
% lane mask and appended by stable compressed stores, already in chronological
% lane order.  The uncommon path is entered only after a hash-level prefilter and
% a raw-record check establish a relevant equal-hash, different-identity case.
% For the first complete window, the special search is for records matching the
% previously saved warm-up hash; for an expiry, it searches for records matching
% the newly reduced minimum hash.  The kernel then constructs event slots for
% all four lanes, rearranges the interleaved slot-major data into one row per
% chronological window, and stably compresses away only maximum sentinels.  It
% does not deduplicate two genuine records.

% \begin{figure}[H]
% \centering
% \begin{tikzpicture}[font=\footnotesize,>=Latex]
%   \node[rewrite box,fill=rewriteBlue!12,minimum width=3.00cm] (mins)
%     at (1.80,4.55) {four adjacent-window minima\\in chronological lanes};
%   \node[rewrite box,minimum width=3.00cm] (events)
%     at (5.40,4.55) {event tests\\entering $\leq$ old or old expired};
%   \draw[rewrite flow] (mins)--(events);

%   \node[rewrite out,align=center,minimum width=3.30cm] (common) at (9.10,5.35)
%     {ordinary lane mask\\two stable compressed stores};
%   \node[rewrite min,align=center,minimum width=3.30cm] (rare) at (9.10,3.62)
%     {raw-distinct equal-hash tie\\build event-slot matrix};
%   \draw[rewrite flow] (events.east) -- ++(.45,0) |- (common.west);
%   \draw[rewrite flow] (events.east) -- ++(.45,0) |- (rare.west);

%   \node[rewrite box,minimum width=3.30cm] (ordered) at (13.00,3.62)
%     {rearrange to one row per window\\stable removal of sentinels};
%   \draw[rewrite flow] (rare)--(ordered);

%   \node[rewrite note,text width=3.80cm,anchor=north] at (9.10,2.72)
%     {This matrix restores event order. It is not the A2 transpose used to put
%      far-apart chunk traces into rows, and it is not a deduplication pass.};
%   \node[rewrite out,align=center,minimum width=3.00cm] (flush) at (1.80,2.05)
%     {end of input\\flush final saved record};
%   \draw[rewrite flow] (mins.south) -- (flush.north);
% \end{tikzpicture}
% \caption{Output processing after the adjacent-window reduction.  The common
% path directly compresses chronological event lanes.  The rare matrix path is
% required only to preserve the order and multiplicity of raw-distinct tied
% records.  Neither path removes equal hashes merely because their values are
% equal.}
% \label{fig:rewrite-exact-event-path}
% \end{figure}

% For DNA5 in the generic chronological implementation, an ambiguous base
% resets the exact $k$-mer context and advances an invalid maximum candidate
% through the minimizer state.  It does not trigger an artificial
% end-of-sequence flush.  Fragment boundaries likewise preserve the continuation
% state.  These rules are outside the DNA4 B8 fast path, but they explain why the
% chronological route owns explicit continuation and finalization operations.

% \subsection{Algorithmic differences}

% \begin{table}[H]
% \centering
% \small
% \begin{tabularx}{\linewidth}{p{.19\linewidth}p{.18\linewidth}p{.20\linewidth}X}
% \toprule
% Configured method & Meaning of one SIMD lane & Hash construction & Sliding-minimum work \\
% \midrule
% \texttt{simd\_minimap2\_}\newline\texttt{reference} &
% One of several separated input chunks &
% Exact canonical code and minimap-style hashing & Saved minimum; a lane expiry
% causes a conditional shared scan of all $w$ ring slots \\
% \texttt{bounded\_nthash\_}\newline\texttt{simd} &
% One of exactly eight separated input streams &
% Rolling 32-bit ntHash-like forward and reverse state; score $F+R$ & Blockwise
% prefix and suffix minima; a fixed suffix rebuild every $w$ candidates \\
% \texttt{two\_stack\_}\newline\texttt{nthash}:\newline
% chronological K15/W10\newline fast path & One of up to four active
% adjacent windows inside a group & Up to eight exact canonical 30-bit codes per block,
% then minimap integer permutation & Ten shifted hash-vector comparison steps
% for every group, with corresponding metadata loads; an unconditional
% four-lane reduction of up to four complete windows \\
% \texttt{two\_stack\_}\newline\texttt{nthash}:\newline resident route &
% One of separated resident input
% streams & Select by the high half of canonical-code multiplication plus local
% position; emit $\phi(C\mathbin{\&}\mathord{\sim}1)$ & Prefix and suffix minima
% on a fixed schedule \\
% \bottomrule
% \end{tabularx}
% \caption{The algorithms differ independently along three axes: lane geometry,
% hash construction, and sliding-minimum data structure.  Calling all of them
% ``SIMD minimizers'' does not make their inner loops equivalent.}
% \label{tab:rewrite-backend-source-differences}
% \end{table}

% The adjacent chronological kernel performs $w$ shifted hash comparison steps
% for each $L=4$-window group, i.e. fixed $\Theta(w)$ vector work per group and
% $\Theta(w/L)$ vector reduction steps per candidate.  Each step also moves the
% raw and ordinal metadata selected by the hash comparison.  Its
% window-reduction loop has a branch-stable, value-independent schedule and
% reuses each shifted hash load across four neighboring lanes; the subsequent
% event path still contains data-dependent tie branches.  The A2 bounded kernel
% instead pays an amortized constant number of minimum operations per candidate
% by rebuilding $w$ suffix aggregates once per $w$ steps.  A1 retains the best
% scalar common case but sometimes pays a complete ring scan, and SIMD lockstep
% makes one lane's expiry activate that work for the whole vector.  These are
% three distinct algorithms, not three names for the same rescan.

% \paragraph{Implementation anchors.}
% The dispatch distinction is visible in
% \texttt{bioseq/backend/pipeline/runtime\_contract.cppm}, lines 81--103, and
% \texttt{minimizer\_scalar\_highway.cppm}, lines 455--466 and 914--1018.
% The bounded ntHash table, rolling recurrence, two-stack update, tie endpoints,
% and eight-stream geometry are in \texttt{bounded\_nthash\_kernel.cppm}, lines
% 45--49, 317--376, 516--607, 618--849, 904--1005, 1178--1217, and
% 1299--1401.  B8 admission is in
% \texttt{two\_stack\_nthash/exact\_vectors.cppm}, lines 49--80, while the
% nonfragmented and fragmented driver choices are in
% \texttt{two\_stack\_nthash/drivers.cppm}, lines 120--129, 319--328, and
% 523--746.  The chronological B8 hashing and its 24-byte scratch are in
% \texttt{two\_stack\_nthash/exact\_k15\_w10.cppm}, lines 595--766; its shifted
% window reduction and event paths are at lines 38--87 and 154--487.  The exact
% minimap permutation is in \texttt{hashers/minimap\_int.cppm}, lines 85--101,
% generic ambiguity handling and event logic are in
% \texttt{two\_stack\_nthash/exact\_chronological.cppm}, lines 312--568 and
% 571--753, and finalization is in
% \texttt{two\_stack\_nthash/continuation.cppm}, lines 494--599.  The resident
% selection and emission scores are in
% \texttt{two\_stack\_nthash/resident\_stream.cppm}, lines 223--451 and
% 584--631.



% \section{TRASHLOLZ}





% Past V8 adopted the paper's useful way of finding window minima: each lane
% followed a far-apart input piece, every lane kept prefix and suffix summaries,
% and the suffix pass ran on a fixed schedule.  The complete algorithm was still
% different.  V8 selected with a cheaper key derived from the canonical packed
% $k$-mer, deferred the full minimap-style scramble until output, resolved an
% equal key with its packed local position, and appended selected lanes through
% a scalar loop.  It did not implement the paper's output transpose and
% cross-stream join removal.  Its ring and summaries were fixed arrays; no
% \texttt{std::map} was part of the minimum calculation.

% The inspected V8 snapshot also contains two representation errors, separate
% from the intended algorithm.  For the usual 32-bit, $k=15$ form, the mask used
% for the multiplied ordering key retains only fourteen useful hash bits.  In a
% 64-bit instantiation of the same expression, its two masks do not overlap and
% the ordering hash becomes zero.  The code also stores strand in the low bit of
% the packed canonical word and clears that bit before calculating the emitted
% minimap-style hash.  A word whose original low bit was one therefore receives
% the wrong emitted hash.  These errors are additional reasons not to treat V8
% output as an exact minimap2 stream; they are not properties of the
% prefix/suffix method itself.

% \begin{figure}[H]
% \centering
% \begin{tikzpicture}[font=\footnotesize]
%   \node[anchor=west,font=\bfseries] at (0,8.10)
%     {(a) One V8 lane follows one far-apart, overlapping input piece};
%   \foreach \r/\x in {0/.65,1/2.25,2/3.85,3/5.45} {
%     \pgfmathsetmacro{\yy}{7.38-.52*\r}
%     \node[anchor=east] at (.48,\yy) {lane \r};
%     \draw[fill=rewriteGray!20] (\x,\yy-.22) rectangle (\x+.72,\yy+.22);
%     \draw[fill=rewriteBlue!14,draw=rewriteBlue!90!black]
%       (\x+.72,\yy-.22) rectangle (\x+2.32,\yy+.22);
%   }
%   \node[rewrite note,text width=4.30cm,anchor=west] at (8.05,6.62)
%     {Neighboring lanes read some of the same bases. Their minimum state does
%      not cross the lane boundary.};

%   \node[anchor=west,font=\bfseries] at (0,5.12)
%     {(b) SIMD calculates four independent window minima};
%   \foreach \r/\lab in {0/L0,1/L1,2/L2,3/L3}
%     \node[anchor=east] at (1.05,4.43-.48*\r) {\lab};
%   \foreach \r/\c/\v in {
%       0/0/9,0/1/4,0/2/7,0/3/2,
%       1/0/8,1/1/3,1/2/6,1/3/5,
%       2/0/6,2/1/1,2/2/8,2/3/4,
%       3/0/7,3/1/5,3/2/2,3/3/6}
%     \node[rewrite cell] at (1.48+.67*\c,4.43-.48*\r) {\v};
%   \foreach \r/\c/\v in {0/3/2,1/1/3,2/1/1,3/2/2}
%     \node[rewrite min] at (1.48+.67*\c,4.43-.48*\r) {\v};
%   \node[rewrite box,minimum width=2.35cm,text width=1.95cm] (fixed)
%     at (5.60,3.72) {fixed prefix/suffix combine};
%   \draw[rewrite flow] (4.02,3.72)--(fixed.west);
%   \foreach \r/\v in {0/2,1/3,2/1,3/2}
%     \node[rewrite min] at (7.34,4.43-.48*\r) {\v};
%   \node[rewrite note,text width=4.15cm,anchor=west] at (7.95,3.72)
%     {The minimum update is vectorized. The branch is controlled by the shared
%      step counter, not by a lane's hash.};

%   \node[anchor=west,font=\bfseries] at (0,2.02)
%     {(c) Sparse output returns to scalar lane handling};
%   \node[rewrite box,minimum width=2.30cm,text width=1.95cm] (store)
%     at (2.00,.74) {store hash, position, strand, and mask};
%   \node[rewrite box,minimum width=2.25cm,text width=1.90cm] (loop)
%     at (5.00,.74) {scalar loop over lane 0 to lane 3};
%   \node[rewrite out,minimum width=1.65cm] at (8.00,1.09) {lane 1 record};
%   \node[rewrite out,minimum width=1.65cm] at (8.00,.37) {lane 3 record};
%   \draw[rewrite flow] (store)--(loop);
%   \draw[rewrite flow] (loop.east)-|(7.05,1.09);
%   \draw[rewrite flow] (loop.east)-|(7.05,.37);
%   \node[rewrite note,text width=3.6cm,anchor=west] at (9.10,.74)
%     {Only repetition within one lane is suppressed. Overlap records from two
%      neighboring lanes can remain.};
% \end{tikzpicture}
% \caption{Past V8 lane processing. Input lanes are far apart, the
% prefix/suffix minimum is calculated with SIMD operations, and sparse records
% are appended by a scalar lane loop. V8 has an input rearrangement for distant
% pieces, but it has no SIMD-Minimizers output rearrangement and join removal.}
% \label{fig:rewrite-v8-lanes}
% \end{figure}

% \begin{table}[H]
% \centering
% \footnotesize
% \begin{tabularx}{\linewidth}{p{.18\linewidth}XX}
% \toprule
% Question & SIMD-Minimizers & Past V8 \\
% \midrule
% What is one lane? & One of eight far-apart, overlapping input streams & One
% far-apart, overlapping input piece \\
% How is the hash advanced? & Rolling forward and reverse ntHash & Rolling
% canonical packed word with a cheaper multiplied ordering key; full
% minimap-style mix delayed until output \\
% How is the minimum found? & Block prefixes and suffixes & The same
% prefix/suffix minimum structure \\
% How is an equal minimum chosen? & For an odd complete span, more G/T than A/C
% selects the leftmost occurrence; otherwise the rightmost occurrence is
% selected & Packed local position selects one earlier equal key;
% complete minimap2 tie events are not reproduced \\
% How is output collected? & Rearrange vector-step output into stream rows and
% remove repeated join identities & Store vectors, inspect lanes in scalar code,
% and suppress repetition only within a lane \\
% Is a tree map used? & No & No \\
% \bottomrule
% \end{tabularx}
% \caption{The shared prefix/suffix minimum does not make V8 and SIMD-Minimizers the
% same complete algorithm. Their hash, equal-minimum rule, and output handling
% differ.}
% \label{tab:rewrite-sea-v8-differences}
% \end{table}

% \section{Current BioSeq algorithms and the route selected by the benchmark}
% \label{sec:rewrite-current-bioseq}

% The public name
% \begin{center}
% \footnotesize
% \texttt{bioseq\_block\_simd\_dp\_}\\
% \texttt{two\_stack\_nthash\_no\_index}
% \end{center}
% preserves names from earlier development.  Static routing is more specific:
% the named large no-index case maps target pieces and the query through the
% state-carrying adjacent-window SIMD routine.  Its active window-minimum loop
% does not call the separate whole-input prefix/suffix routine and does not call
% the bounded rolling-ntHash routine.  The phrase \emph{adjacent-window
% routine} is used below because lane~0 handles one window, lane~1 handles the
% very next window, and so on.

% \begin{figure}[H]
% \centering
% \begin{tikzpicture}[font=\footnotesize]
%   \node[rewrite box,minimum width=7.25cm,fill=rewriteGray!15] (name)
%     at (6.30,6.35)
%     {named E2E row:\\
%      \texttt{bioseq\_block\_simd\_dp\_}\\
%      \texttt{two\_stack\_nthash\_no\_index}};
%   \node[rewrite box,minimum width=4.20cm] (select) at (6.30,5.05)
%     {public selector \texttt{TWO\_STACK\_NTHASH}\\chooses the exact SIMD family};
%   \draw[rewrite flow] (name)--(select);

%   \node[rewrite box,minimum width=4.25cm,fill=rewriteBlue!12] (target)
%     at (3.20,3.45) {each target segment\\state carried to the next segment};
%   \node[rewrite box,minimum width=4.25cm,fill=rewriteBlue!12] (query)
%     at (9.40,3.45) {query sketch\\same exact seed rules};
%   \draw[rewrite flow] (select.south) |- (target.north);
%   \draw[rewrite flow] (select.south) |- (query.north);

%   \node[rewrite box,minimum width=5.15cm,fill=rewriteOrange!14] (adj)
%     at (6.30,1.90) {adjacent-window SIMD block\\exact minimap-style hash and ties};
%   \draw[rewrite flow] (target.south) |- (adj.west);
%   \draw[rewrite flow] (query.south) |- (adj.east);

%   \node[rewrite note,text width=4.45cm,anchor=west,fill=rewriteGray!10]
%     at (.35,.35) {Not selected here: whole-input \texttt{vector\_stream}
%     two-stack control.};
%   \node[rewrite note,text width=4.45cm,anchor=west,fill=rewriteGray!10]
%     at (7.85,.35) {Not selected here: bounded rolling-ntHash scheme.};
% \end{tikzpicture}
% \caption{What the named no-index E2E row actually runs. Both target segments
% and the query use the state-carrying adjacent-window SIMD routine. The
% words ``two stack'' and ``nthash'' in the public row name do not describe the
% active inner minimum or hash in this route.}
% \label{fig:rewrite-current-e2e-route}
% \end{figure}

% \subsection{The adjacent-window calculation}

% Let $B$ be the number of candidate positions handled together.  The routine
% retains the last $k-1$ bases and the last $w-1$ candidate records.  It adds $B$
% new consecutive bases, forms $B$ consecutive $k$-mers, calculates the exact
% minimap-style hash for each, and places the retained and new candidates next
% to one another.  It then makes $w$ shifted SIMD loads.  Lane $j$ is the whole
% window beginning at the $j$th candidate, so lanes are neighboring windows in
% this routine.  This is the opposite of V8, where lanes are far-apart input
% pieces.

% \begin{figure}[H]
% \centering
% \begin{tikzpicture}[font=\footnotesize]
%   \node[anchor=west,font=\bfseries] at (0,9.05)
%     {(a) Retained context plus $B$ new consecutive positions};
%   \node[rewrite box,minimum width=4.0cm,fill=rewriteGray!18] (basesold)
%     at (2.35,8.28) {last $k-1$ bases};
%   \node[rewrite box,minimum width=2.40cm,fill=rewriteBlue!14] (basesnew)
%     at (5.75,8.28) {$B$ new bases};
%   \node[rewrite box,minimum width=4.05cm] (make) at (9.65,8.28)
%     {$B$ consecutive canonical $k$-mers\\and exact minimap-style hashes};
%   \draw[rewrite flow] (basesold.east)--(basesnew.west);
%   \draw[rewrite flow] (basesnew.east)--(make.west);

%   \node[anchor=west,font=\bfseries] at (0,7.28)
%     {(b) Four neighboring windows are calculated together ($B=4,w=4$ shown)};
%   \node[anchor=east] at (1.78,6.58) {candidate row};
%   \foreach \i/\v/\sty in {0/7/rewrite old,1/6/rewrite old,2/5/rewrite old,
%                             3/9/rewrite new,4/4/rewrite new,5/8/rewrite new,
%                             6/3/rewrite new}
%     \node[\sty] at (2.20+.70*\i,6.58) {\v};

%   \foreach \r/\lab in {0/$W_0$,1/$W_1$,2/$W_2$,3/$W_3$}
%     \node[anchor=east] at (1.78,5.82-.54*\r) {\lab};
%   \foreach \r/\c/\v in {
%       0/0/7,0/1/6,0/2/5,0/3/9,
%       1/0/6,1/1/5,1/2/9,1/3/4,
%       2/0/5,2/1/9,2/2/4,2/3/8,
%       3/0/9,3/1/4,3/2/8,3/3/3}
%     \node[rewrite cell] at (2.20+.70*\c,5.82-.54*\r) {\v};
%   \foreach \r/\c/\v in {0/2/5,1/3/4,2/2/4,3/3/3}
%     \node[rewrite min] at (2.20+.70*\c,5.82-.54*\r) {\v};

%   \node[rewrite box,minimum width=2.45cm] (reduce) at (6.38,5.02)
%     {$w$ shifted loads\\and $\leq$ reduction};
%   \draw[rewrite flow] (5.02,5.02)--(reduce.west);
%   \foreach \r/\v in {0/5,1/4,2/4,3/3}
%     \node[rewrite min] at (8.32,5.82-.54*\r) {\v};
%   \node[rewrite note,text width=3.80cm,anchor=west] at (8.92,5.02)
%     {One lane is one complete window. The next lane is the window one
%      candidate position to the right.};

%   \node[anchor=west,font=\bfseries] at (0,3.15)
%     {(c) The ordinary output path writes only changed selections};
%   \node[anchor=east] at (1.93,2.45) {previous};
%   \node[anchor=east] at (1.93,1.88) {current};
%   \node[anchor=east] at (1.93,1.31) {write?};
%   \foreach \i/\v in {0/5@3,1/5@3,2/4@5,3/4@5}
%     \node[rewrite old,minimum width=.93cm] at (2.48+.98*\i,2.45) {\v};
%   \foreach \i/\v in {0/5@3,1/4@5,2/4@5,3/3@7}
%     \node[rewrite new,minimum width=.93cm] at (2.48+.98*\i,1.88) {\v};
%   \foreach \i/\v in {0/0,1/1,2/0,3/1}
%     \node[rewrite cell,minimum width=.93cm] at (2.48+.98*\i,1.31) {\v};
%   \node[rewrite box,minimum width=2.35cm,fill=rewriteGreen!13] (compress)
%     at (7.45,1.88) {compressed store\\of marked lanes};
%   \draw[rewrite flow] (5.95,1.88)--(compress.west);
%   \node[rewrite out,minimum width=1.28cm] at (9.55,2.20) {5@3};
%   \node[rewrite out,minimum width=1.28cm] at (9.55,1.55) {4@5};
%   \node[rewrite note,text width=2.55cm,anchor=west] at (10.35,1.88)
%     {3@7 remains saved until a later event.};
% \end{tikzpicture}
% \caption{The current adjacent-window minimum with small values for
% legibility. Gray candidates come from the preceding block and blue candidates
% are new. Four lanes hold four neighboring windows. The real $k=15,w=10$
% specialization reads ten shifted vectors and produces $B$ window minima.}
% \label{fig:rewrite-current-adjacent-windows}
% \end{figure}

% The comparison uses $\leq$, so a later equal-hash position replaces an earlier
% one.  Ordinary changes form one lane mask and are written in a compressed
% batch.  Equal hashes at several different positions can require several
% records for one window.  That rare case rearranges an event matrix so all
% records for $W_0$ precede all records for $W_1$, and so on.  The rearrangement
% restores output order; it is not a deduplication pass.

% \begin{figure}[H]
% \centering
% \begin{tikzpicture}[font=\footnotesize]
%   \node[font=\bfseries] at (2.45,5.32) {before: one row per event slot};
%   \foreach \c/\lab in {0/$W_0$,1/$W_1$,2/$W_2$,3/$W_3$}
%     \node at (1.35+.78*\c,4.72) {\lab};
%   \foreach \r/\lab in {0/old,1/tie 1,2/tie 2}
%     \node[anchor=east] at (.87,4.12-.56*\r) {\lab};
%   \foreach \r/\c/\v in {0/0/a,0/1/b,0/2/--,0/3/f,
%                           1/0/c,1/1/--,1/2/e,1/3/g,
%                           2/0/d,2/1/--,2/2/--,2/3/--}
%     \node[rewrite cell] at (1.35+.78*\c,4.12-.56*\r) {\v};

%   \node[rewrite box,minimum width=1.85cm] (swap) at (5.60,3.55)
%     {rows become\\columns};
%   \draw[rewrite flow] (4.28,3.55)--(swap.west);
%   \draw[rewrite flow] (swap.east)--(6.83,3.55);

%   \node[font=\bfseries] at (9.18,5.32) {after: one row per window};
%   \foreach \c/\lab in {0/old,1/tie 1,2/tie 2}
%     \node at (8.05+.90*\c,4.72) {\lab};
%   \foreach \r/\lab in {0/$W_0$,1/$W_1$,2/$W_2$,3/$W_3$}
%     \node[anchor=east] at (7.53,4.12-.56*\r) {\lab};
%   \foreach \r/\c/\v in {0/0/a,0/1/c,0/2/d,
%                           1/0/b,1/1/--,1/2/--,
%                           2/0/--,2/1/e,2/2/--,
%                           3/0/f,3/1/g,3/2/--}
%     \node[rewrite cell] at (8.05+.90*\c,4.12-.56*\r) {\v};

%   \node[rewrite out,minimum width=2.25cm] at (2.10,.85) {a\quad c\quad d};
%   \node[rewrite out,minimum width=.85cm] at (3.82,.85) {b};
%   \node[rewrite out,minimum width=.85cm] at (4.78,.85) {e};
%   \node[rewrite out,minimum width=1.55cm] at (6.02,.85) {f\quad g};
%   \node[rewrite note,text width=5.25cm,anchor=west] at (7.05,.85)
%     {Every named record remains. Only empty sentinel cells disappear.};
% \end{tikzpicture}
% \caption{Rare equal-minimum output in the current adjacent-window routine.
% The matrix is rearranged to put records in window order. A compressed store
% removes empty cells; it does not compare or remove two real records.}
% \label{fig:rewrite-current-event-order}
% \end{figure}

% \subsection{State across input fragments}

% A normal storage-fragment boundary does not end the biological sequence.  The
% routine retains at most $k-1$ bases, fewer than one incomplete vector group,
% $w-1$ candidates, and the selected record of the last complete window.  An
% ambiguous base is different: it clears the valid $k$-mer run and advances a
% maximum-valued candidate through the window.

% \begin{figure}[H]
% \centering
% \begin{tikzpicture}[font=\footnotesize]
%   \draw[fill=rewriteGray!18] (.55,4.55) rectangle (5.35,5.25);
%   \node at (2.95,4.90) {packed fragment $f_i$};
%   \draw[fill=rewriteBlue!14,draw=rewriteBlue!90!black]
%     (7.05,4.55) rectangle (11.85,5.25);
%   \node at (9.45,4.90) {packed fragment $f_{i+1}$};
%   \draw[rewrite flow] (5.35,4.90)--(7.05,4.90);
%   \node[above,align=center] at (6.20,5.05) {storage boundary\\no reset};

%   \node[rewrite box,minimum width=11.30cm,fill=rewriteGray!12] (state)
%     at (6.20,3.02) {retained state: $k-1$ bases \quad $w-1$ candidates
%       \quad incomplete group \quad selected record};
%   \draw[rewrite flow] (4.25,4.50)|-(state.north);
%   \draw[rewrite flow] (state.north)-|(8.15,4.50);

%   \node[rewrite box,minimum width=9.90cm,fill=rewriteOrange!13] (amb)
%     at (6.20,1.40) {ambiguous base: clear the valid $k$-mer context and
%       insert a maximum-valued candidate};
%   \draw[rewrite flow] (state)--(amb);
%   \node[rewrite note,text width=10.8cm,anchor=west] at (.80,.20)
%     {Only end-of-input performs the final saved-record write. A fragment
%      boundary and an ambiguous base therefore have different effects.};
% \end{tikzpicture}
% \caption{Bounded continuation in the current adjacent-window SIMD path. State is carried
% across a normal storage boundary, so a window can cross it without rereading
% an overlapping chunk. An ambiguous base resets the valid $k$-mer run but does
% not masquerade as end-of-input.}
% \label{fig:rewrite-current-continuation}
% \end{figure}

% \begin{algorithm}[H]
% \caption{Current $k=15,w=10$ adjacent-window block, simplified}
% \label{alg:rewrite-current-adjacent-block}
% \small
% \begin{enumerate}
%   \item Join retained base context with $B$ new positions and build $B$
%         consecutive forward and reverse-complement $k$-mers.
%   \item Choose each canonical $k$-mer, mark invalid or palindromic candidates
%         with a maximum sentinel, and calculate the exact minimap-style hash.
%   \item Join the retained $w-1$ candidate tail with the $B$ new candidates.
%         Read $w$ shifted vectors and reduce with $\leq$; lane $j$ now contains
%         the minimum of the $j$th neighboring window.
%   \item Form masks for a smaller entering value, an expired saved record,
%         first-window output, and final output.
%   \item On the common path, compressed-store the marked records.  If equal
%         hashes at different identities are present, build the tie-event
%         matrix, rearrange it into window order, and compressed-store all real
%         cells.
%   \item Retain the final base context, candidate tail, incomplete group, and
%         selected record for the next input fragment.
% \end{enumerate}
% \end{algorithm}

% \subsection{Why this complete route can be faster than V8}

% The current minimum operation is not an asymptotically cheaper form of V8's
% two-stack minimum.  V8 rebuilds $w$ suffix entries once every $w$ steps.  The
% current routine makes $w$ shifted comparisons for every group of $B$ adjacent
% windows and also calculates the exact selection hash.  A speed advantage, when
% measured, must therefore come from the complete route around that reduction:

% \begin{itemize}
%   \item new $k$-mers come from one short consecutive input region, rather than
%         from an input rearrangement of far-apart pieces;
%   \item state crosses ordinary fragments directly, without overlapping lane
%         reads and their repeated join records;
%   \item the ordinary output is written as one compressed batch rather than by
%         visiting every lane in scalar code; and
%   \item the larger exact tie calculation runs only after a hash precheck says
%         it is needed.
% \end{itemize}

% Those are code-level reasons, not timing results.  Section
% \ref{sec:rewrite-performance-evidence} keeps the older measurements separate
% and does not present them as a V8-versus-current measurement.

% \subsection{Other current minimizer routes}

% BioSeq has several routes because they do not all produce the same seed
% stream. In particular, the three values assigned to
% \texttt{dispatch.params.minimizer\_algorithm} select three different ways of
% making minimizers. They do not select the prefilter, chaining method, or final
% alignment method; those are chosen by other fields. The same minimizer setting
% is used for the query and target so that a query seed can be found in the
% target table.

% \subsubsection{What the three ordinary benchmark settings mean}

% The ordinary alignment benchmark normally supplies DNA4 packed into 32-bit
% words, with the complete query and target already available in memory, and
% uses $k=15,w=10$. Under those conditions the three assignments in the
% benchmark source mean the following.

% \begin{description}
%   \item[\texttt{two\_stack\_nthash}.]
%     The four \texttt{bioseq\_prefilter\_*} rows and the six ordinary alignment
%     rows whose names contain \texttt{two\_stack\_nthash} run the whole-input
%     descendant of V8. One lane follows one far-apart part of the sequence. It
%     advances packed forward and reverse-complement $k$-mers, updates a prefix
%     minimum, and rebuilds suffix minima once every $w$ positions. Long inputs
%     advance two independent Highway vector groups. Despite the setting's
%     historical name, this route does not calculate ntHash. The E2E caller
%     shown earlier uses the same setting but supplies the target as pieces and
%     requires exact minimap-style output for its target table, so it
%     deliberately selects the adjacent-window algorithm instead.

%   \item[\texttt{bounded\_nthash\_simd}.]
%     The two ordinary rows whose names contain \texttt{bounded\_nthash} call
%     the bounded-input entry point. The query is processed as one range and the
%     target as one or more ranges. Every range uses eight logical ntHash
%     streams and the prefix/suffix minimum. This is the maintained successor of
%     the old \texttt{PAPER\_NTHASH\_SIMD} setting. It is the BioSeq algorithm
%     closest to SIMD-Minimizers, but its exact hash and tied-minimum output differ from
%     the paper, as explained below.

%   \item[\texttt{simd\_minimap2\_reference}.]
%     This setting is used by the ordinary slow-reference row. It runs the naive
%     SIMD rescan from
%     Figure~\ref{fig:rewrite-naive-simd-rescan}. One lane follows one far-apart
%     sequence part and saves one minimum. If any lane loses its saved minimum,
%     the shared loop reads all $w$ vector slots; masks ensure that only expired
%     lanes change their answer.
% \end{description}

% Words such as \texttt{ksw2}, \texttt{block}, \texttt{prefilter}, or
% \texttt{vulkan\_tiled\_dp} elsewhere in an ordinary row name describe later
% mapping work. They do not change the minimizer route just described. Highway
% chooses the available processor target and physical lane count; there is no
% separate handwritten AVX algorithm. If a required input shape is unsupported,
% the checked SIMD setting reports failure. It does not silently change to the
% scalar minimizer after entering that checked route.

% \subsubsection{The whole-input prefix/suffix route}

% When \texttt{two\_stack\_nthash} receives one complete, already available
% DNA4 sequence stored in 32-bit words with $k=15,w=10$, it can choose the
% whole-input descendant of V8.  This is the route reached by the ordinary
% in-memory benchmark names that contain \texttt{two\_stack\_nthash}.  Each lane
% owns a far-apart sequence piece.  The code shifts packed forward and
% reverse-complement words, ranks their canonical word with the multiplied
% ordering key, updates a prefix minimum, and rebuilds suffix minima every $w$
% steps.  Long inputs advance two independent Highway lane groups together.

% This route still does not calculate ntHash.  It applies the minimap-style mix
% only when it emits, stores SIMD results to arrays, and then visits output lanes
% with scalar code.  It has no state format for continuing through ordinary
% input fragments.  For these reasons it cannot be substituted for the
% adjacent-window route while a target index and query are being built from
% pieces.

% \begin{table}[H]
% \centering
% \footnotesize
% \begin{tabularx}{\linewidth}{p{.22\linewidth}p{.22\linewidth}X
%   p{.10\linewidth}}
% \toprule
% Minimizer setting & Input for this row & Actual algorithm & Named
% no-index E2E? \\
% \midrule
% \texttt{scalar} & DNA4 or DNA5 scalar input & saved minimum plus occasional
% ring rescan & no \\
% \shortstack[l]{\texttt{simd\_minimap2\_}\\\texttt{reference}} &
% one complete DNA4 input & distant
% lane streams with one saved minimum per lane and a shared SIMD rescan & no \\
% \texttt{two\_stack\_nthash} & a target table or input supplied in pieces &
% neighboring windows, full $w$-slot SIMD reduction, minimap-style hash & yes \\
% \texttt{two\_stack\_nthash} & complete DNA4 stored in 32-bit words,
% $k=15,w=10$ &
% distant lane streams with V8-descended prefix/suffix minima & no \\
% \texttt{bounded\_nthash\_simd} & DNA4 query or bounded target range,
% $k=15,w=10$ & eight rolling-ntHash streams with prefix/suffix minima & no \\
% \shortstack[l]{\texttt{vulkan\_minimap\_}\\\texttt{integer}} &
% supported GPU and input & many GPU
% tiles, each using a serial minimap-style ring/rescan & no \\
% \bottomrule
% \end{tabularx}
% \caption{Current minimizer settings and their actual inner algorithms.  The
% \texttt{two\_stack\_nthash} selector is the confusing case: its caller and
% input shape choose between two different algorithms.}
% \label{tab:rewrite-current-routing}
% \end{table}

% \begin{table}[H]
% \centering
% \footnotesize
% \begin{tabularx}{\linewidth}{p{.40\linewidth}XX}
% \toprule
% Configured run name or family & Query minimizers & Target minimizers \\
% \midrule
% \shortstack[l]{Named E2E \texttt{two\_stack\_nthash\_}\\
% \texttt{no\_index}} & adjacent-window SIMD &
% adjacent-window SIMD while building a temporary target table \\
% \shortstack[l]{Named E2E \texttt{two\_stack\_nthash\_}\\
% \texttt{just\_index}} & adjacent-window SIMD &
% matching adjacent-window target index built before mapping \\
% Ordinary \texttt{bioseq\_prefilter\_*} and names containing
% \texttt{two\_stack\_nthash} & whole-input V8 descendant & whole-input V8
% descendant \\
% Ordinary \texttt{simd\_minimap2\_slow} & direct SIMD rescan & direct SIMD
% rescan \\
% Ordinary names containing \texttt{bounded\_nthash} & bounded rolling ntHash &
% bounded rolling ntHash \\
% Scalar BioSeq rows & scalar rescan & scalar rescan \\
% Vulkan-minimizer rows & Vulkan tile rescan & Vulkan tile rescan or its
% device target-index builder \\
% External minimap2 rows & bundled minimap2 sketch & bundled minimap2 sketch or
% its prebuilt index \\
% \bottomrule
% \end{tabularx}
% \caption{What the configured test families actually execute.  This mapping
% comes from the current call paths, not from interpreting a historical word in
% the displayed run name.  No current configured run calls the old V8 source
% body directly.}
% \label{tab:rewrite-configured-routes}
% \end{table}

% \begin{table}[H]
% \centering
% \footnotesize
% \begin{tabularx}{\linewidth}{p{.19\linewidth}XX}
% \toprule
% Question & Past V8 & Current adjacent-window E2E path \\
% \midrule
% What is one lane? & One far-apart, overlapping input piece & One of the next
% $B$ neighboring windows \\
% How is input prepared? & Rearrange packed input so every lane receives the
% next base of its own piece & Read consecutive bases and retain $k-1$ bases
% across input pieces \\
% How is the minimum found? & Update a prefix; rebuild suffixes once per $w$
% steps; compare one prefix and suffix & Read $w$ shifted vectors and reduce
% each of $B$ complete windows \\
% Which hash selects it? & Cheaper multiplied ordering key; full minimap-style
% mix delayed until output & Exact minimap-style hash before selection \\
% How are normal records written? & Store vectors and inspect lanes in scalar
% code & One event mask and one compressed SIMD store \\
% What is rearranged? & Input values for far-apart pieces & Rare equal-minimum
% output records, solely to restore window order \\
% Is there a deduplication pass? & No; suppression is within one lane & No;
% unchanged selections are suppressed before output and real ties remain \\
% \bottomrule
% \end{tabularx}
% \caption{Past V8 and the current route selected by the named no-index E2E
% row. The current code did not replace V8 with a cheaper two-stack minimum; it
% changed lane meaning, input handling, exactness, and record writing.}
% \label{tab:rewrite-v8-current}
% \end{table}

% \subsubsection{How BioSeq bounded ntHash differs from SIMD-Minimizers}

% Earlier versions called this route \texttt{PAPER\_NTHASH\_SIMD}, and some
% development text abbreviated its output format as PNT1.  It was not removed.
% The production route was renamed \texttt{bounded\_nthash\_simd}; its numeric
% selector stayed unchanged. Target tables made by this algorithm still carry
% the four-letter marker PNT1, which tells the reader that the bounded hash and
% tied-position rules were used. The current bounded kernel is therefore the
% maintained successor of that earlier production kernel. A different
% experimental helper that accepted an unpacked array of base numbers was
% removed, but the query, index, and benchmark routes never called that helper.
% Restoring the old production source would create a second copy of an
% algorithm that is already present.

% Here \emph{bounded} means that one call receives a finite DNA range and clear
% limits saying which window starts belong to that range. The kernel rereads
% enough neighboring bases to rebuild its state; it does not carry rolling-hash
% or prefix/suffix state from one storage fragment to the next.

% The algorithm always has eight logical input pieces. A Highway target with
% eight physical 32-bit lanes holds them in one vector group; a target with four
% physical lanes holds them in two groups. The checked implementation rejects
% other physical lane counts instead of changing the algorithm.

% The current bounded algorithm works as follows.

% \begin{enumerate}
%   \item It divides one DNA4 range into eight far-apart pieces. Each piece is
%         assigned its own interval of window starts; these intervals do not
%         overlap. To finish windows at a boundary, a piece needs $k+w-2$ bases
%         to its right. Packed reads are rounded to whole 32-bit words and can
%         therefore include more bases.
%   \item It rearranges packed input words so the next SIMD value contains the
%         next DNA letter from every far-apart piece. This is the same kind of
%         \emph{input} rearrangement used by SIMD-Minimizers. It removes no output.
%   \item Every lane maintains a forward and reverse-complement rolling hash.
%         The two unsigned 32-bit values are added, with overflow wrapping back
%         into 32 bits, to obtain the value compared inside a window.
%   \item Every lane updates a prefix minimum on each step and rebuilds $w$
%         suffix minima once per $w$ candidates. One prefix/suffix comparison
%         gives the next window minimum. This is the paper's two-stack idea.
%   \item If the minimum hash occurs once, the algorithm writes that occurrence.
%         If it occurs several times, it writes the leftmost and rightmost
%         occurrences; it does not write equal occurrences between them.
%   \item While one piece is processed, a record is skipped when its hash,
%         position, and forward/reverse mark are identical to the
%         preceding selected event for that piece. The first record of each
%         later piece is held back and compared with the last record of the
%         preceding non-empty piece. An identical boundary record is written
%         once.
% \end{enumerate}

% Step 6 is important: bounded ntHash does \emph{not} implement the paper's
% output-matrix rearrangement. SIMD-Minimizers first produces one selected position for
% every window, rearranges a block of those positions into one ordered row per
% input piece, and then compares neighboring positions in SIMD. Bounded ntHash
% instead compares an event with the preceding event while each lane advances,
% stores the lane results in small arrays, and appends them with a scalar loop.
% Both methods avoid many repeated outputs from neighboring windows, but they do
% so at different times and with different instructions.

% The hash and tied-minimum rules also differ. These are not small storage
% details: they can select different positions from the same DNA sequence.

% \begin{table}[H]
% \centering
% \footnotesize
% \begin{tabularx}{\linewidth}{p{.18\linewidth}XX}
% \toprule
% Question & SIMD-Minimizers & BioSeq bounded ntHash \\
% \midrule
% How is the rolling value updated? & Use one-bit rotations and a fixed lookup
% table. The packed DNA order is A, C, T, G. & Use seven-bit rotations and
% BioSeq's own four lookup values. BioSeq's packed DNA order is A, C, G, T. The
% resulting 32-bit numbers are therefore not the same as the paper's numbers. \\
% How much of the hash chooses the minimum? & Compare only the upper 16 bits;
% the lower 16 bits of the packed value hold a local position. & Compare all 32
% hash bits and keep the position separately. Two values tied by the paper's
% upper 16 bits may be ordered differently here. \\
% What happens when the minimum occurs twice? & For odd $k+w-1$, count G/T
% against A/C over the complete span. More G/T chooses the leftmost minimum;
% otherwise choose the rightmost. Return one occurrence. & Return both extreme
% occurrences: leftmost first and rightmost second. \\
% How is a forward/reverse value recorded? & The whole-span count chooses an
% endpoint; the paper does not attach a per-$k$-mer strand bit to that returned
% position. & Mark an endpoint as reverse only when its reverse rolling value
% is smaller than its forward value. Equality is marked forward. \\
% How are repeated positions removed? & Rearrange an output matrix into one row
% per input piece, compare neighboring positions, compact the changes, and
% remove the repeated position at a piece join. & Compare with the preceding
% record during each lane's walk and compare the last/first records at a piece
% join. There is no output-matrix rearrangement. \\
% \bottomrule
% \end{tabularx}
% \caption{A direct comparison of SIMD-Minimizers and BioSeq bounded ntHash. They share
% far-apart input pieces, an input rearrangement, rolling forward/reverse hash
% state, and prefix/suffix minima. They calculate different hash numbers, make
% different choices for tied minima, and remove repeated output by different
% methods.}
% \label{tab:rewrite-sea-bounded-nthash}
% \end{table}

% For BioSeq's usual $k=15,w=10$, the complete span has
% $k+w-1=24$ bases, which is even. This does not prevent rolling ntHash or the
% prefix/suffix minimum, but it does not satisfy the odd-span condition of the
% paper's canonical endpoint rule. BioSeq bounded ntHash does not use that
% rule; it reports both extreme occurrences of a tied minimum.

% \begin{figure}[H]
% \centering
% \begin{tikzpicture}[font=\footnotesize]
%   \node[anchor=west,font=\bfseries] at (0,7.85)
%     {(a) One window has the same minimum at two positions};
%   \node[anchor=east] at (2.15,6.25) {position};
%   \node[anchor=east,align=right] at (2.15,5.58) {compared\\key};
%   \foreach \i in {0,...,3}
%     \node[rewrite cell] at (2.65+.78*\i,6.25) {\i};
%   \foreach \i/\v in {0/9,1/2,2/5,3/2}
%     \node[rewrite cell] at (2.65+.78*\i,5.58) {\v};
%   \node[rewrite min] at (2.65+.78*1,5.58) {2};
%   \node[rewrite min] at (2.65+.78*3,5.58) {2};

%   \node[rewrite box,text width=3.30cm,minimum width=3.65cm,
%         fill=rewriteBlue!10] at (7.25,6.12)
%     {SIMD-Minimizers, odd span:\\more G/T selects position 1; otherwise
%      position 3; write one};
%   \node[rewrite box,text width=3.30cm,minimum width=3.65cm,
%         fill=rewriteOrange!12] at (11.15,6.12)
%     {BioSeq bounded:\\write positions 1 and 3; mark each reverse only if its
%      reverse hash is smaller};

%   \node[anchor=west,font=\bfseries] at (0,4.35)
%     {(b) Repeated selections from neighboring windows};
%   \node[rewrite box,minimum width=2.25cm] (events) at (1.70,3.05)
%     {positions\\$[1,1,5,5]$};
%   \node[rewrite box,text width=3.70cm,minimum width=4.05cm,
%         fill=rewriteBlue!10] (seaout) at (5.40,3.05)
%     {SIMD-Minimizers:\\rearrange output into one row per input piece; compare
%      neighboring positions};
%   \node[rewrite out,minimum width=1.85cm] (searesult) at (8.70,3.05)
%     {$[1,5]$};
%   \draw[rewrite flow] (events)--(seaout);
%   \draw[rewrite flow] (seaout)--(searesult);

%   \node[rewrite box,minimum width=2.25cm] (eventsb) at (1.70,1.05)
%     {positions\\$[1,1,5,5]$};
%   \node[rewrite box,text width=3.70cm,minimum width=4.05cm,
%         fill=rewriteOrange!10] (boundout) at (5.40,1.05)
%     {bounded:\\compare each event with the preceding event while its input
%      piece is being processed};
%   \node[rewrite out,minimum width=1.85cm] (boundresult) at (8.70,1.05)
%     {$[1,5]$};
%   \draw[rewrite flow] (eventsb)--(boundout);
%   \draw[rewrite flow] (boundout)--(boundresult);
%   \node[rewrite note,text width=3.05cm,anchor=west] at (9.85,2.05)
%     {At an input-piece join, bounded ntHash also compares the preceding
%      piece's last record with the next piece's first record.};
% \end{tikzpicture}
% \caption{Two visible output differences. The small numbers in panel (a) are
% the keys seen by the minimum comparison: the paper obtains them from the
% upper 16 hash bits, while bounded ntHash normally compares all 32 bits.
% SIMD-Minimizers
% returns one end of this tie, while bounded ntHash returns both ends. Panel (b)
% shows that the same simple run of repeated positions can reduce to the same
% two outputs, but SIMD-Minimizers first rearranges an output matrix whereas bounded ntHash
% compares successive events during the lane walk.}
% \label{fig:rewrite-sea-bounded-output}
% \end{figure}

% Because these choices can produce different hashes and different positions,
% the bounded query must be compared with a target table made by the bounded
% algorithm. The named \texttt{two\_stack\_nthash} E2E rows do not select it.
% The ordinary rows whose names contain \texttt{bounded\_nthash} do.

% \subsubsection{Vulkan: many serial tile walks run together}

% On the GPU, a segment is a memory-sized part of a large input, not one
% minimizer window.  A segment includes context around an owned core.  Inside
% the uploaded segment, each shader invocation owns a smaller tile.  It starts
% before the tile, replays enough bases to rebuild rolling and window state, and
% then writes only events triggered inside its half-open owned interval.  The
% window minimum inside one tile is the same saved-minimum plus ring-rescan
% method shown in Figure~\ref{fig:rewrite-minimap-rescan}; it is not a
% prefix/suffix minimum.

% Because the number of output events is not known in advance, the compact
% route runs a count pass, an exclusive scan of tile counts, and an emit pass.
% Each tile receives a non-overlapping output range before it writes.

% \begin{figure}[H]
% \centering
% \begin{tikzpicture}[font=\footnotesize]
%   \node[anchor=west,font=\bfseries] at (0,10.45)
%     {(a) One memory segment has context around its owned core};
%   \draw[fill=rewriteGray!20] (.55,9.50) rectangle (2.55,10.18);
%   \draw[fill=rewriteBlue!14,draw=rewriteBlue!90!black]
%     (2.55,9.50) rectangle (10.55,10.18);
%   \draw[fill=rewriteGray!20] (10.55,9.50) rectangle (12.55,10.18);
%   \node at (1.55,9.84) {context};
%   \node at (6.55,9.84) {owned core};
%   \node at (11.55,9.84) {context};

%   \node[anchor=west,font=\bfseries] at (0,8.88)
%     {(b) Tile invocations own adjacent trigger intervals};
%   \draw[fill=rewriteGray!20] (.55,7.90) rectangle (2.40,8.58);
%   \draw[fill=rewriteBlue!14,draw=rewriteBlue!90!black]
%     (2.40,7.90) rectangle (6.15,8.58);
%   \draw[fill=rewriteGray!20] (4.30,6.80) rectangle (6.15,7.48);
%   \draw[fill=rewriteBlue!14,draw=rewriteBlue!90!black]
%     (6.15,6.80) rectangle (9.90,7.48);
%   \node at (1.48,8.24) {replay};
%   \node at (4.28,8.24) {tile $i$ owns $[b_i,e_i)$};
%   \node at (5.23,7.14) {replay};
%   \node at (8.03,7.14) {tile $i+1$ owns next interval};
%   \draw[densely dashed,rewriteGray!90] (6.15,6.48)--(6.15,8.85);
%   \node[anchor=west] at (10.20,7.69)
%     {$e_i=b_{i+1}$};

%   \node[anchor=west,font=\bfseries] at (0,5.98)
%     {(c) One tile walks positions in order and uses a ring with rescan};
%   \node at (.65,4.98) {$\cdots$};
%   \foreach \i/\v/\sty in {0/2/rewrite min,1/7/rewrite old,2/6/rewrite old,
%                             3/5/rewrite new,4/9/rewrite cell,5/4/rewrite cell,
%                             6/3/rewrite cell}
%     \node[\sty] at (1.25+.72*\i,4.98) {\v};
%   \draw[very thick,rewriteGray!90] (3.06,4.55)--(3.06,5.41);
%   \draw[draw=rewriteOrange!90!black,very thick]
%     (.88,4.50) rectangle (3.82,5.46);
%   \node[rewrite note,text width=5.15cm,anchor=west] at (7.05,4.55)
%     {Roll one canonical $k$-mer, push it into the $w$-slot ring, compare it
%      with the saved minimum, and rescan only when that saved slot leaves.};
%   \node[anchor=west,rewriteGray!90!black] at (1.10,4.15)
%     {events during replay are not written};

%   \node[anchor=west,font=\bfseries] at (0,3.35)
%     {(d) Compact output allocation};
%   \node[rewrite box,minimum width=2.35cm] (count) at (1.65,2.45)
%     {count events\\per tile};
%   \node[rewrite box,minimum width=2.35cm] (scan) at (4.72,2.45)
%     {exclusive scan\\of counts};
%   \node[rewrite box,minimum width=2.35cm] (emit) at (7.79,2.45)
%     {replay tile and\\write events};
%   \node[rewrite box,minimum width=2.65cm,fill=rewriteGreen!14] (ranges)
%     at (11.05,2.45) {separate output\\range for each tile};
%   \draw[rewrite flow] (count)--(scan);
%   \draw[rewrite flow] (scan)--(emit);
%   \draw[rewrite flow] (emit)--(ranges);
%   \node[rewrite note,text width=11.3cm,anchor=west] at (.65,.78)
%     {Parallelism comes from many tile invocations. Positions inside one tile
%      are processed in order because the rolling $k$-mer and saved minimum
%      depend on earlier positions.};
% \end{tikzpicture}
% \caption{How the current Vulkan sketch divides work. Context is read again,
% but only the owned core or tile can publish an event. Each tile uses a serial
% ring/rescan state machine. Count--scan--emit gives tiles exact, separate
% output ranges while many tiles execute at the same time.}
% \label{fig:rewrite-vulkan-sketch}
% \end{figure}

% \subsubsection{Why target and query must use the same minimizer rules}

% A target table is looked up by the hash written by its minimizer algorithm.
% BioSeq bounded ntHash and the minimap-style algorithms can write different
% hashes and different tied positions for the same DNA. If the target table was
% made by bounded ntHash but the query was made by adjacent-window SIMD, many
% query hashes would simply not exist in that table even when the DNA matched.

% BioSeq therefore first asks which minimizer rules made the target table. It
% then chooses an implementation that produces exactly those rules. If that
% implementation cannot handle the alphabet, $k$, $w$, input layout, or
% available Highway lane count, the operation reports an error. Secretly
% switching to another minimizer algorithm would produce a plausible-looking
% but incorrect lookup.

% \begin{figure}[H]
% \centering
% \begin{tikzpicture}[font=\footnotesize]
%   \node[rewrite box,minimum width=3.20cm,fill=rewriteBlue!12] (scheme)
%     at (6.20,6.75) {How was the target table made?};
%   \node[rewrite box,minimum width=3.10cm] (pntproof) at (2.70,5.00)
%     {Can bounded ntHash\\handle this input?};
%   \node[rewrite box,minimum width=3.30cm] (mmproof) at (9.70,5.00)
%     {Which minimap-style\\implementation was requested?};
%   \coordinate (routefork) at (6.20,5.90);
%   \draw[thick] (scheme.south)--(routefork);
%   \draw[rewrite flow] (routefork)-|(pntproof.north)
%     node[pos=.38,above] {bounded ntHash};
%   \draw[rewrite flow] (routefork)-|(mmproof.north)
%     node[pos=.38,above] {minimap-style};

%   \node[rewrite box,minimum width=2.75cm,fill=rewriteGreen!14] (pnt)
%     at (2.70,3.45) {bounded ntHash\\SIMD};
%   \node[rewrite box,minimum width=2.75cm,fill=rewriteOrange!14] (rejectp)
%     at (2.70,1.90) {reject mismatch};
%   \draw[rewrite flow] (pntproof)--(pnt) node[midway,right] {yes};
%   \draw[rewrite flow] (pntproof.west)-|(rejectp.west)
%     node[pos=.25,left] {no};

%   \node[rewrite box,minimum width=2.35cm,fill=rewriteGreen!14] (exact)
%     at (7.10,3.45) {adjacent-window\\SIMD};
%   \node[rewrite box,minimum width=2.35cm,fill=rewriteGreen!14] (vk)
%     at (9.70,3.45) {Vulkan};
%   \node[rewrite box,minimum width=2.35cm,fill=rewriteGreen!14] (scalar)
%     at (12.30,3.45) {Scalar};
%   \draw[rewrite flow] (mmproof.south)-|(exact.north);
%   \draw[rewrite flow] (mmproof)--(vk);
%   \draw[rewrite flow] (mmproof.south)-|(scalar.north);

%   \node[rewrite box,minimum width=5.25cm,fill=rewriteOrange!14] (fail)
%     at (9.70,1.55) {requested implementation cannot handle input: error};
%   \draw[rewrite flow] (exact.south)--++(0,-.45)
%     -| ($(fail.north)+(-1.65,0)$);
%   \draw[rewrite flow] (vk)--(fail);
%   \draw[rewrite flow] (scalar.south)--++(0,-.45)
%     -| ($(fail.north)+(1.65,0)$);
% \end{tikzpicture}
% \caption{Why a checked implementation cannot silently change minimizer
% algorithms. The query must use the same hash and tied-position rules as the
% target table. Bounded ntHash therefore cannot fall through to a minimap-style
% implementation, and a failed adjacent-window or Vulkan check is reported.}
% \label{fig:rewrite-backend-selection}
% \end{figure}

% \section{Performance evidence}
% \label{sec:rewrite-performance-evidence}

% Timing claims are grouped here so they cannot be mistaken for algorithm
% definitions.  The first four charts preserve an older three-run measurement
% series.  They are useful history, but they predate the current adjacent-window
% route and use a different binary and experiment scope.  Their source retained
% means but not the raw samples or sample standard deviations; the figures
% therefore show no error bars.  They must not be used as a direct measured
% V8-versus-current comparison.

% \subsection{Older end-to-end series}

% \begin{figure}[H]
%   \centering
% \begin{CacheMeCode}{python, custom include command={\input{\robExtAddCachePathAndName{\robExtFinalHash.pgf}}}, set placeholder eval={__LINEWIDTH__}{\lenToCmNoUnit[in]{\linewidth}}}
%     import matplotlib
%     matplotlib.use("pgf")
%     import matplotlib.pyplot as plt
%     import numpy as np
%     matplotlib.rcParams.update({
%       "font.family": "serif",
%       "font.serif": [],
%       "text.usetex": True,
%     })
%     fig, ax = plt.subplots()
%     fig.set_size_inches(__LINEWIDTH__, 0.55*__LINEWIDTH__, forward=True)
%     buckets = ['47 kbp', '49 kbp', '52 kbp', '903 kbp', '5.0 Mbp']
%     minimap2 = [20.11, 20.10, 20.13, 190.95, 1060.13]
%     scalar = [14.56, 14.94, 15.09, 188.93, 1171.0]
%     simd_v6 = [15.21, 15.90, 16.05, 201.68, 1207.99]
%     v8 = [14.64, 15.07, 15.24, 183.77, 1131.87]
%     historical_v13 = [13.78, 14.02, 13.82, 172.38, 1113.3]
%     x = np.arange(len(buckets))
%     width = 0.16
%     ax.bar(x-2*width, minimap2, width, label='minimap2', color='#6F6F6F')
%     ax.bar(x-width, scalar, width, label='scalar', color='#B8B8B8')
%     ax.bar(x, simd_v6, width, label='naive SIMD rescan', color='#4C78A8')
%     ax.bar(x+width, v8, width, label='past V8', color='#F28E2B')
%     ax.bar(x+2*width, historical_v13, width,
%            label='historical V12/V13 route', color='#59A14F')
%     ax.set_xticks(x)
%     ax.set_xticklabels(buckets)
%     ax.set_yscale('log')
%     ax.set_ylabel('wall time per pair (ms, log scale)', fontsize=8)
%     ax.tick_params(axis='both', labelsize=7)
%     ax.legend(fontsize=7, frameon=False, loc='upper left', ncol=2)
%     ax.grid(True, alpha=0.25, axis='y', which='both')
%     print(get_filename_from_extension('.pgf'))
%     plt.savefig(get_filename_from_extension('.pgf'), bbox_inches='tight')
% \end{CacheMeCode}
%   \caption{Older end-to-end wall times across five sequence lengths. Each bar
%   is the recorded mean of three runs; raw samples and sample standard
%   deviations were not retained with this chart. ``Historical V12/V13 route''
%   is not the current adjacent-window E2E route. The log scale allows the short
%   and long inputs to share one axis.}
%   \label{fig:rewrite-historical-headline}
% \end{figure}

% \subsection{How the older minimizer phase changed}

% \begin{figure}[H]
%   \centering
% \begin{CacheMeCode}{python, custom include command={\input{\robExtAddCachePathAndName{\robExtFinalHash.pgf}}}, set placeholder eval={__LINEWIDTH__}{\lenToCmNoUnit[in]{\linewidth}}}
%     import matplotlib
%     matplotlib.use("pgf")
%     import matplotlib.pyplot as plt
%     import numpy as np
%     matplotlib.rcParams.update({
%       "font.family": "serif",
%       "font.serif": [],
%       "text.usetex": True,
%     })
%     fig, ax = plt.subplots()
%     fig.set_size_inches(__LINEWIDTH__, 0.53*__LINEWIDTH__, forward=True)
%     names = ['V6\nrescan', 'V8\nprefix/suffix', 'V8.5\ndelayed mix',
%              'V11\ntwo streams', 'V12\ninlined', 'scalar']
%     inner = [23.88, 19.90, 13.60, 13.60, 9.23, 0.0]
%     wrapper = [14.94, 12.43, 8.62, 8.62, 5.76, 12.50]
%     x = np.arange(len(names))
%     ax.bar(x, wrapper, label='recorded wrapper field', color='#4C78A8')
%     ax.bar(x, inner, bottom=wrapper, label='recorded inner field',
%            color='#F28E2B')
%     ax.axhline(12.50, linestyle='--', linewidth=1, color='#6F6F6F',
%                label='scalar recorded field')
%     for i, total in enumerate(np.array(inner)+np.array(wrapper)):
%         ax.text(i, total+0.5, f'{total:.1f}', ha='center', fontsize=7)
%     ax.set_xticks(x)
%     ax.set_xticklabels(names)
%     ax.set_ylabel('recorded time per 903-kbp pair (ms)', fontsize=8)
%     ax.tick_params(axis='both', labelsize=7)
%     ax.legend(fontsize=7, frameon=False, loc='upper right')
%     ax.grid(True, alpha=0.25, axis='y')
%     print(get_filename_from_extension('.pgf'))
%     plt.savefig(get_filename_from_extension('.pgf'), bbox_inches='tight')
% \end{CacheMeCode}
%   \caption{Older minimizer-phase fields for one 903-kbp case, recorded as
%   means of three runs without retained raw samples or sample standard
%   deviations. The chart preserves the original stacked presentation. Before
%   using the stack as an additive time, the instrumentation must confirm that
%   the wrapper and inner fields are exclusive rather than nested.}
%   \label{fig:rewrite-historical-phase-evolution}
% \end{figure}

% \subsection{Older complete-run phase composition}

% \begin{figure}[H]
%   \centering
% \begin{CacheMeCode}{python, custom include command={\input{\robExtAddCachePathAndName{\robExtFinalHash.pgf}}}, set placeholder eval={__LINEWIDTH__}{\lenToCmNoUnit[in]{\linewidth}}}
%     import matplotlib
%     matplotlib.use("pgf")
%     import matplotlib.pyplot as plt
%     import numpy as np
%     matplotlib.rcParams.update({
%       "font.family": "serif",
%       "font.serif": [],
%       "text.usetex": True,
%     })
%     fig, ax = plt.subplots()
%     fig.set_size_inches(__LINEWIDTH__, 0.58*__LINEWIDTH__, forward=True)
%     configs = ['scalar', 'V6 rescan', 'V8', 'V12 before V13',
%                'V12 after V13']
%     phases = [
%       ('sketch fields', [12.50,38.82,32.43,14.99,14.52], '#4C78A8',''),
%       ('target index', [16.0,14.3,14.6,26.38,23.83], '#A0A0A0',''),
%       ('seed hits', [6.5,11.0,11.3,12.32,12.0], '#59A14F',''),
%       ('chain', [16.0,8.9,8.9,8.88,8.78], '#F28E2B',''),
%       ('error estimate', [0.7,7.2,7.1,7.88,0.85], '#F28E2B','//'),
%       ('alignment', [87.5,86.4,86.5,87.25,85.42], '#6F6F6F','..'),
%     ]
%     bottom = np.zeros(len(configs))
%     for label, values, color, hatch in phases:
%         ax.bar(configs, values, bottom=bottom, label=label, color=color,
%                hatch=hatch, edgecolor='#404040', linewidth=.25)
%         bottom += np.array(values)
%     ax.set_ylabel('recorded time per 903-kbp pair (ms)', fontsize=8)
%     ax.tick_params(axis='both', labelsize=7)
%     ax.legend(fontsize=6.5, frameon=False, loc='upper right', ncol=2)
%     ax.grid(True, alpha=0.25, axis='y')
%     plt.xticks(rotation=13, ha='right')
%     print(get_filename_from_extension('.pgf'))
%     plt.savefig(get_filename_from_extension('.pgf'), bbox_inches='tight')
% \end{CacheMeCode}
%   \caption{Older phase composition for one 903-kbp experiment. Values are
%   recorded means of three runs; no sample standard deviations are available.
%   The chart separates a minimizer-side change from the later error-estimation
%   change, but it does not describe the current named no-index route.}
%   \label{fig:rewrite-historical-phase-stack}
% \end{figure}

% \subsection{Older percentage view of the same end-to-end series}

% \begin{figure}[H]
%   \centering
% \begin{CacheMeCode}{python, custom include command={\input{\robExtAddCachePathAndName{\robExtFinalHash.pgf}}}, set placeholder eval={__LINEWIDTH__}{\lenToCmNoUnit[in]{\linewidth}}}
%     import matplotlib
%     matplotlib.use("pgf")
%     import matplotlib.pyplot as plt
%     import numpy as np
%     matplotlib.rcParams.update({
%       "font.family": "serif",
%       "font.serif": [],
%       "text.usetex": True,
%     })
%     fig, ax = plt.subplots()
%     fig.set_size_inches(__LINEWIDTH__, 0.50*__LINEWIDTH__, forward=True)
%     buckets = ['47 kbp', '49 kbp', '52 kbp', '903 kbp', '5.0 Mbp']
%     scalar = np.array([14.56,14.94,15.09,188.93,1171.0])
%     v6 = np.array([15.21,15.90,16.05,201.68,1207.99])
%     v8 = np.array([14.64,15.07,15.24,183.77,1131.87])
%     old13 = np.array([13.78,14.02,13.82,172.38,1113.3])
%     x = np.arange(len(buckets))
%     width = .27
%     ax.bar(x-width, (v6-scalar)/scalar*100, width,
%            label='naive SIMD rescan', color='#4C78A8')
%     ax.bar(x, (v8-scalar)/scalar*100, width,
%            label='past V8', color='#F28E2B')
%     ax.bar(x+width, (old13-scalar)/scalar*100, width,
%            label='historical V12/V13 route', color='#59A14F')
%     ax.axhline(0, linewidth=.7, color='#6F6F6F')
%     ax.axhline(-10, linestyle='--', linewidth=1, color='#A0A0A0',
%                label='historical -10\% target')
%     ax.set_xticks(x)
%     ax.set_xticklabels(buckets)
%     ax.set_ylabel('difference from scalar (\%, negative is faster)', fontsize=8)
%     ax.tick_params(axis='both', labelsize=7)
%     ax.legend(fontsize=7, frameon=False, loc='upper right', ncol=2)
%     ax.grid(True, alpha=0.25, axis='y')
%     print(get_filename_from_extension('.pgf'))
%     plt.savefig(get_filename_from_extension('.pgf'), bbox_inches='tight')
% \end{CacheMeCode}
%   \caption{Percentage form of the older three-run end-to-end means. This is a
%   transformation of Figure~\ref{fig:rewrite-historical-headline}, not an
%   independent experiment. Raw samples and sample standard deviations were not
%   retained with the chart.}
%   \label{fig:rewrite-historical-speedup}
% \end{figure}

% \subsection{A later historical CPU checkpoint}

% A later experiment recorded four complete observations per configuration and
% reported their median.  Its bounded-ntHash row used a target index built with
% that scheme, while the minimap2 and adjacent-window rows used their matching
% seed schemes.  It is therefore
% a comparison of complete valid configurations, not interchangeable function
% bodies.

% \begin{figure}[H]
% \centering
% \begin{tikzpicture}[x=.82cm,y=.80cm,font=\footnotesize]
%   \draw[rewrite flow] (0,0)--(10.9,0) node[right] {sketch time (s)};
%   \foreach \x in {0,2,4,6,8,10} {
%     \draw (\x,-.08)--(\x,.08);
%     \node[below] at (\x,-.08) {\x};
%   }
%   \fill[rewriteGreen!75] (0,.48) rectangle (2.923,1.10);
%   \node[anchor=east] at (-.20,.79) {bounded ntHash};
%   \node[anchor=west] at (3.05,.79) {2.923};
%   \fill[rewriteGray!75] (0,1.50) rectangle (4.733,2.12);
%   \node[anchor=east] at (-.20,1.81) {minimap2};
%   \node[anchor=west] at (4.86,1.81) {4.733};
%   \fill[rewriteBlue!75] (0,2.52) rectangle (9.256,3.14);
%   \node[anchor=east] at (-.20,2.83) {adjacent SIMD};
%   \node[anchor=west] at (9.39,2.83) {9.256};
% \end{tikzpicture}
% \caption{Sketch component in a later historical CPU checkpoint. Each value is
% the median of four observations. The retained summary does not provide a
% sample standard deviation. Bounded ntHash uses a different target/query seed scheme, so
% the chart compares complete configurations rather than identical output.}
% \label{fig:rewrite-historical-cpu-checkpoint}
% \end{figure}

% \subsection{GPU before and after the rolling, all-tiles rewrite}

% \begin{table}[H]
% \centering
% \small
% \begin{tabular}{lrrr}
% \toprule
% GPU sketch version & per-call sketch & throughput & complete probe \\
% \midrule
% rebuild-heavy, too few active lanes & 272.0\,ms & 73\,Mbp/s & 21.83\,s \\
% rolling update, all-tile schedule & 40.3\,ms & 494\,Mbp/s & 11.70\,s \\
% \bottomrule
% \end{tabular}
% \caption{One controlled GPU before/after measurement. The two rows use the
% same recorded workload and exact DNA4 output check. Repetition count and
% sample spread were not retained in the source summary, so these values should
% not be given error bars.}
% \label{tab:rewrite-gpu-progression}
% \end{table}

% \subsection{Limit of the retained measurements}

% None of the preserved charts runs past V8 and the current adjacent-window
% kernel with the same input, output rule, binary, and timed boundary.  The
% figures therefore do not establish a numerical speedup between those two
% algorithms.  The supported conclusion is the code-level explanation in
% Section~\ref{sec:rewrite-current-bioseq}: the current route removes V8's
% far-apart input preparation and scalar lane-output loop, but its window
% minimum itself performs more comparisons.

% \subsection{Related seed designs retained as visual reference}

% The following two figures are retained because they explain terms used when
% comparing sketching work, although they are not steps of the E2E minimizer
% route described above.  A spaced seed masks positions in one span.  A linked
% seed chooses separated components.  BLEND constructs a fuzzy fingerprint by
% letting token hashes vote on each output bit.

% \begin{figure}[H]
% \centering
% \begin{tikzpicture}[font=\footnotesize]
%   \node[anchor=west,font=\bfseries] at (0,8.85) {(a) Minimizer};
%   \foreach \i/\txt/\rank/\sty in {
%       0/GCTATTA/161/rewrite cell,1/CTATTAC/2/rewrite min,
%       2/TATTACC/252/rewrite cell,3/ATTACCT/69/rewrite cell}
%     \node[\sty,minimum width=2.25cm,minimum height=.82cm,align=center]
%       at (1.35+2.55*\i,8.02) {\texttt{\txt}\\rank \rank};
%   \node[rewrite box,minimum width=2.50cm] (mselect) at (6.45,6.72)
%     {choose the\\lowest rank};
%   \node[rewrite out,minimum width=2.50cm] (mout) at (10.15,6.72)
%     {\texttt{CTATTAC}};
%   \foreach \x in {1.35,3.90,6.45,9.00}
%     \draw[rewrite flow] (\x,7.58)|-(mselect.north);
%   \draw[rewrite flow] (mselect)--(mout);

%   \node[anchor=west,font=\bfseries] at (0,5.72) {(b) Spaced seed};
%   \node[rewrite box,minimum width=2.25cm] (sin) at (1.50,4.82)
%     {\texttt{GCTATTA}\\\texttt{GATACTA}};
%   \node[rewrite box,minimum width=2.55cm,fill=rewriteBlue!12] (mask)
%     at (4.65,4.82) {mask \texttt{1011011}\\keep the 1 positions};
%   \node[rewrite out,minimum width=2.25cm,align=center] (skeep)
%     at (7.80,4.82)
%     {\texttt{GTATA}\\\texttt{GTATA}};
%   \node[rewrite box,minimum width=2.15cm] (shash) at (10.65,4.82)
%     {hash retained\\symbols};
%   \draw[rewrite flow] (sin)--(mask);
%   \draw[rewrite flow] (mask)--(skeep);
%   \draw[rewrite flow] (skeep)--(shash);

%   \node[anchor=west,font=\bfseries] at (0,3.55) {(c) Linked seed};
%   \node[rewrite out,minimum width=1.65cm,align=center] (l1) at (1.30,2.52)
%     {$s_1$\\\texttt{GCTA}};
%   \node[rewrite box,minimum width=2.25cm] (lw2) at (3.72,2.52)
%     {$W_2$ candidates\\\texttt{CCTT}, \texttt{TTAA}};
%   \node[rewrite out,minimum width=1.65cm,align=center] (l2) at (6.08,2.52)
%     {choose\\\texttt{TTAA}};
%   \node[rewrite box,minimum width=2.25cm] (lw3) at (8.45,2.52)
%     {$W_3$ candidates\\\texttt{GTGA}, \texttt{TGGA}};
%   \node[rewrite out,minimum width=1.65cm,align=center] (l3) at (10.80,2.52)
%     {choose\\\texttt{TGGA}};
%   \draw[rewrite flow] (l1)--(lw2);
%   \draw[rewrite flow] (lw2)--(l2);
%   \draw[rewrite flow] (l2)--(lw3);
%   \draw[rewrite flow] (lw3)--(l3);
%   \node[rewrite note,text width=10.9cm,anchor=west] at (.55,.72)
%     {A linked key contains the chosen components. Bases between its component
%      windows are not part of the key.};
% \end{tikzpicture}
% \caption{Three distinct seed constructions. A minimizer selects one complete
% $k$-mer by rank. A spaced seed hashes only positions marked by a binary mask.
% A linked seed chooses components from explicit downstream windows and joins
% those components into one key.}
% \label{fig:rewrite-related-seeds}
% \end{figure}

% \begin{figure}[H]
% \centering
% \begin{tikzpicture}[font=\footnotesize]
%   \node[anchor=west,font=\bfseries] at (0,6.75)
%     {(a) BLEND-I on two nearby contiguous seeds};
%   \node[rewrite box,minimum width=2.20cm] (b1) at (1.45,5.66)
%     {\texttt{GCTATTAC}\\\texttt{GCTACTAC}};
%   \node[rewrite box,minimum width=2.50cm,fill=rewriteBlue!12] (tokens)
%     at (4.55,5.66) {overlapping\\tokens};
%   \node[rewrite box,minimum width=2.40cm] (bits) at (7.52,5.66)
%     {ordinary token\\hash bits};
%   \node[rewrite box,minimum width=2.15cm,fill=rewriteOrange!14] (vote)
%     at (10.28,5.66) {signed vote\\for each bit};
%   \node[rewrite out,minimum width=2.05cm,align=center] (fp) at (12.82,5.66)
%     {fuzzy\\fingerprint};
%   \draw[rewrite flow] (b1)--(tokens);
%   \draw[rewrite flow] (tokens)--(bits);
%   \draw[rewrite flow] (bits)--(vote);
%   \draw[rewrite flow] (vote)--(fp);
%   \node[rewrite note,text width=11.9cm,anchor=west] at (.55,4.12)
%     {Similar seeds share many token hashes, so many bit votes remain the same.
%      A shared fingerprint is possible, not guaranteed.};

%   \node[anchor=west,font=\bfseries] at (0,3.25)
%     {(b) BLEND-S on linked components};
%   \node[rewrite box,minimum width=2.75cm] (linked) at (1.75,2.15)
%     {\texttt{GCTA|TTAA|TGGA}};
%   \node[rewrite box,minimum width=2.35cm,fill=rewriteBlue!12] (ctok)
%     at (5.00,2.15) {component\\token hashes};
%   \node[rewrite box,minimum width=2.20cm,fill=rewriteOrange!14] (cvote)
%     at (7.90,2.15) {bit vote};
%   \node[rewrite box,minimum width=2.20cm] (lookup) at (10.70,2.15)
%     {index lookup};
%   \node[rewrite out,minimum width=2.10cm,align=center] (verify)
%     at (13.35,2.15)
%     {candidates\\to verify};
%   \draw[rewrite flow] (linked)--(ctok);
%   \draw[rewrite flow] (ctok)--(cvote);
%   \draw[rewrite flow] (cvote)--(lookup);
%   \draw[rewrite flow] (lookup)--(verify);
%   \node[rewrite note,text width=11.9cm,anchor=west] at (.55,.57)
%     {A fuzzy-key match proposes candidates. It does not establish that the
%      underlying sequence regions are equal.};
% \end{tikzpicture}
% \caption{Fuzzy fingerprint construction in BLEND~\cite{blend}. Token hashes
% vote independently on each output bit. Nearby seeds can retain enough votes
% to share a fingerprint, after which downstream work still verifies the
% candidate match.}
% \label{fig:rewrite-blend-fuzzy-hash}
% \end{figure}

% \section*{Chapter conclusion}

% Minimap2 saves one minimum and scans the complete window only when that record
% leaves. The naive SIMD port keeps the same rule in every lane, so one expiring
% lane makes the complete lane group scan. SIMD-Minimizers avoids that shared surprise
% with rolling ntHash and fixed prefix/suffix summaries. Past V8 keeps the
% prefix/suffix minimum but changes the hash, equal-minimum rule, and output
% collection. The named current no-index E2E row does something else again: it
% calculates exact minimap-style minima for neighboring windows and writes
% ordinary events with a compressed SIMD store. Vulkan uses many serial tile
% walks, each with a ring and rescan. These routes can be compared only after the
% seed rule and the timed work have been stated explicitly.
